% Modernised reading — fonds Alexandre Grothendieck, shelfmark 69
% (pages 1-119, the whole folder).
%
% This is an INTERPRETATION, not a transcription. It re-reads the pages in
% today's notation and vocabulary, states the hypotheses the manuscript leaves
% implicit, and drops the critical apparatus so that the mathematics reads
% straight through. Everything the brackets and underlines carried moves into
% a footnote. It is held to being mathematically correct as it stands; where
% the manuscript is loose or mistaken, the true statement is what appears here
% and the footnote says what the page has. In any dispute of substance the
% transcription governs: whoever cites Grothendieck must cite batch-01.fr.tex
% (pp. 1-20), batch-02.fr.tex (pp. 21-40), batch-03.fr.tex (pp. 41-60),
% batch-04.fr.tex (pp. 61-80), batch-05.fr.tex (pp. 81-100) or
% batch-06.fr.tex (pp. 101-119).
%
% Pass: Opus 5.5 (claude-opus-5-5), 2026-10-10 — first pass, unchecked against the pages by a human.
% The folder was read whole, its six batches together; this is one document
% for the shelfmark. It was made on Opus 5.5 and has not been measured against
% the reference reading of folder 115 (made on Opus 5). The literature named
% in the footnotes is cited from memory and was not checked against the
% sources. Every count, orbit length, scalar product and root expansion
% asserted below about the 27-vertex graph and E_6 was checked by machine on
% the standard model (the 27 classes xi_i, xi'_i, xi_ij in Z^{1,6}): srg
% parameters (27,10,1,5), |W| = 51840 (group generated by the 36 reflections),
% Lemmas 3 and 4 and the Epilogue labelling, the model of page 13, the rule of
% page 18, the induced 4-, 5-, 6-cycles (1080, 2592, 720), independent sets
% (27, 216, 720, 1080, 648, 72), the stabiliser orbits of page 61, the
% bitriangle rule of page 36, the double-sixes of pages 49 and 63, the simple
% roots and the table of page 92, the 54 minuscule weights, page 89's table.
%
% Scope. Page 1 is a cover (pencilled « (34p) »), page 40 a pink separator
% (« (23 p) »), pages 64, 65, 72, 76, 88, 93, 94, 99, 104, 106 blank
% separators; page 83 is a typed course announcement (USTL, DEUG A, 1976,
% « Polyèdres réguliers »), described only as far as its transcription note
% goes. The inventory's « lettre (s.d.) » is not found in the transcriptions.
% From page 95 on, the archivists' pencilled numbers run two ahead of the page
% numbers used here (which follow the facsimile, as the transcriptions do).
%
% Spine. The graph C of the 27 lines on a cubic surface (« graphe » or
% « complexe cubique »), its automorphism group W = W(E_6), and the root
% system E_6 that sits behind it:
%   pp. 2-10   characterisation: 27 vertices, valency 10, every edge in one
%              triangle => C is determined (Lemmas 1-9, Épilogue).
%   pp. 11-13  the tower point < edge < ... < bitriangle, |W| = 51840; an
%              explicit model.
%   pp. 14-15  the algebra behind the double covers: central extensions,
%              Baer product over Psi(V), anisotropic forms, dim V <= 2.
%   p. 16      the general axioms; card S >= 3(nu-1).
%   pp. 17-39  C with a marked configuration (vertex, edge, triangle, balise,
%              sablier, prism, bitriangle): equivalent data, stabiliser as a
%              Weyl group (D_5, B_4, D_4, F_4, B_3, A_3, A_2), the orthogonal
%              roots, the check |W| = 2^7 3^4 5 each time.
%   pp. 41-60  the 27 lines from the blow-up of P^2 in six points; Pic, the
%              form, lambda = -K; roots, reflections, W = Aut C, the theorem
%              « R(C) is E_6 », simple roots, highest root, Aut R = W x {±1},
%              the map delta -> 3 delta - lambda; an intrinsic definition and
%              the list of the 72 bases.
%   pp. 61-80  a census: configurations, n-free sets, bitriangles and their
%              self-complementarity, polygons and polyhedra.
%   pp. 81-98  differences and sums of vertices, relative positions of roots,
%              the 27 in simple-root coordinates, and the search for the sets
%              S, S_0 (minuscule weights) that would give the 27 back from E_6.
%   pp. 100-118 three normal subgroups; three torsors and a correspondence
%              Gamma in P_1 x P_2 x P_3 (a Latin square); its autotopies;
%              bitorsors; and (p. 113) the axiom that returns to pp. 7-8.
%   p. 105, p. 119  asides (E_7 and its 56; S^2 x S^2, idempotents).
%
% Notation fixed for the document (each departure footnoted where it occurs).
%  - C for the cubic graph throughout (the pages: S on pp. 2-10, C from p. 11,
%    Delta on pp. 45-60); C_0 the standard one (his Delta_0). Delta is kept
%    for a bitriangle (his use on pp. 34-39 and 73-75).
%  - « liés » kept: two vertices are linked when the two lines meet. E(x) the
%    neighbours of x. « Base » = six mutually skew lines (one six of a
%    double-six); « bibase » = double-six {b, b'}.
%  - The form is the intersection form: delta.delta = -1, r.r = -2 for roots;
%    E = lambda^perp is negative definite. lambda = 3 eta - xi = -K.
%  - T = {u,v,w} a triangle, E(u)* = E(u) \ {v,w}, sigma = sigma_T,
%    P_u = E(u)*/sigma. Gamma~ the 32 triangles of C \ T, Gamma = Gamma~/sigma.
%    The extension of p. 14 is written Psi~ (his Gamma~ there), to keep
%    Gamma~ for triangles.
%  - The sets of pp. 95-98 are written S, S_0, S', S'_0 in \mathfrak.
%
% Substantive departures, each footnoted where it occurs:
%  p. 2    Lemma 2's x' lies in E(u)*, not E(v)* (false as written: 4).
%  p. 7    Gamma~ ≅ ∐ (not ∏) of the omega_v(x); V* means V \ {0}.
%  p. 14   proposition stated for 2G = 0 (the case used); p_k alpha_i = -id
%          in general.
%  p. 22   theta = eta - xi_6 = lambda - u (page: eta - xi'_6).
%  p. 23   card W = 27.40.2.4! = 51840 (page: « 27.3^3.5 »).
%  p. 33   Aut of the sablier has order 8 (page: Z/4Z).
%  p. 42   K = -lambda (margin: K = -3 lambda).
%  p. 44   3 eta' = 15 eta - 6 xi (page: -5 xi).
%  p. 49   the double-six of r_ij is (xi_j, xi'_j, xi_ik), (xi_i, xi'_i,
%          xi_jk) as the page has (the transcription's note doubts it; the
%          page is right); p. 51's (xi_i, xi'_j, xi_ik) is not a base.
%  p. 56   i(delta).i(delta') = 6 for linked (page: 9); 3R is the set of
%          differences of NON-linked elements; not 72 vectors of norm -12.
%  p. 57   Z/2 generated by -id_E.
%  p. 58   I and I' free supplied; b) « non lié à i ».
%  p. 61   several orbit lists corrected (2-base, ordered 2-base, 3-base
%          ordered, 4-base, 2-star); 2-stars number 216.
%  p. 63   r_0 = r_1+2r_2+3r_3+2r_4+r_5+2r_6 and r_456 = ... + r_6 (page:
%          3(r_1+r_2+r_3)+...); « -1 si i = i' ».
%  p. 67   B''_5 (216), not B''_4.
%  p. 82   sums of linked vertices have square 0, not 2.
%  p. 86   the parity rule reconciled only on one parity class.
%  p. 90/92 the -xi~_i rows: p. 92 right, p. 90 not; the -xi~'_i rows of
%          p. 92 have the sign of 3 r_0 reversed.
%  pp. 95-98 S defined with the coroot normalisation (struck on the page);
%          the facts for B_r, C_r corrected; « S'_0 = S_0 ∪ R » read as
%          S' = S ∪ R.
%  p. 100  the products are direct; N_1 is abelian.
%
% Links the pages do not write (ours, said so in the body): the axioms of
% p. 2 are those of the generalised quadrangle GQ(2,4); the double covers are
% governed by Q_8 (p. 14 with V of dim 2); W_t of order 1152 is W(F_4); the
% 120 bitriangles are Steiner's trihedral pairs and the self-complementarity
% of p. 73 is the Paley graph of order 9 / the affine plane AG(2,3); the 54
% vectors ±(delta - lambda/3) are the minuscule weights of E_6, the set S of
% p. 95, which answers p. 57's question; pp. 107-118 are the theory of
% autotopies of a group's Latin square; p. 113's axiom is Lemmas 3-4 of p. 3.
% Folder 75 (p. 21) has the hyperoctahedral « structures cubiques »; folder 86
% has the course notes « Polyèdres réguliers » that page 83 announces.
%
% Announced and never established: Lemma on bases p. 20 (uniqueness, in a
% cancelled block); the « autodualité » of p. 38 (taken up pp. 73-75, left
% at the definition of a « dualité »); programme a)-c) of p. 39 (only a)
% done); the intrinsic recovery of Delta~ from (E, R) asked on pp. 56-57
% (pp. 95-98 do not reach E_6); condition (A) of p. 109; the Lemma of p. 80.
\documentclass[11pt,a4paper]{article}
\input{../preamble/grothendieck}

\folder{69}
\pages{1}{119}
\dating{[à partir de 1976]}
\watermark{Édition de démonstration\\interprétation personnelle de l'œuvre}
\foldertitle{Graphes cubiques : notes manuscrites (s.d.), lettre (s.d.) — lecture modernisée du dossier entier}

\begin{document}

\section*{Résumé}

\begin{resume}
Une surface cubique --- l'ensemble des points de l'espace où s'annule un
polynôme de degré trois en trois variables --- contient, si elle est lisse,
exactement vingt-sept droites. C'est un fait du XIX\textsuperscript{e}
siècle, et il en traîne un autre : ces droites se rencontrent selon un dessin
très rigide. Chacune en rencontre dix autres, et deux droites qui se
rencontrent en ont toujours exactement une troisième qui rencontre les deux
(les trois sont dans un même plan). Le dossier part de ce dessin, qu'il
appelle le \emph{graphe cubique} : vingt-sept sommets, deux sommets étant
« liés » quand les droites se coupent.

La première question est celle d'un mathématicien qui veut se passer de la
géométrie : ces trois propriétés --- vingt-sept sommets, dix voisins
chacun, un unique troisième sommet sur chaque arête --- suffisent-elles à
reconstituer tout le dessin ? Il le démontre sur dix feuillets serrés, en
choisissant un triangle et en regardant comment le reste du graphe se répartit
autour de lui. Puis il demande ce qu'on obtient quand on fixe une
configuration du graphe --- un sommet, une arête, un triangle, deux triangles
--- et la réponse revient toujours sous la même forme : la donnée équivaut à
quelque chose d'élémentaire (un ensemble à cinq éléments, un ensemble à six
éléments partagé en deux moitiés\ldots), et le groupe des symétries qui la
respectent est un \emph{groupe de Weyl}, le groupe des symétries d'un système
de vecteurs très régulier. À chaque fois il vérifie que le nombre de
configurations, multiplié par la taille de ce groupe, redonne
$51\,840$, l'ordre du groupe de toutes les symétries du graphe.

Ce nombre est celui du groupe de Weyl de $E_6$, l'un des systèmes de racines
« exceptionnels ». La deuxième moitié du dossier le montre directement : en
réalisant la surface comme un plan où l'on a fait éclater six points, les
vingt-sept droites deviennent des vecteurs d'un réseau, leurs différences
forment les soixante-douze racines de $E_6$, et les symétries du graphe sont
exactement les réflexions qu'elles engendrent. On pense à un cristal : le
graphe est la trace visible, la géométrie des racines est la structure qui la
porte. Il pose alors la question inverse --- peut-on retrouver les
vingt-sept droites à partir des seules racines ? --- et la poursuit sur
plusieurs feuillets, système de racines par système de racines, sans la
conclure ; la réponse est dans ce qu'on appelle aujourd'hui les poids
minuscules.

Les derniers feuillets changent d'échelle. Autour d'un triangle, le graphe se
décrit par trois ensembles et une relation entre eux, une table à double
entrée où chaque ligne et chaque colonne contient chaque valeur une fois --- un
carré latin. Il se demande quand une telle table vient d'un groupe, et quelles
sont ses symétries, d'abord pour un groupe commutatif, puis en général avec des
« bitorseurs ». On voit là une manière de faire : une configuration finie et
concrète est déshabillée jusqu'à ce qu'il n'en reste qu'une structure
algébrique, et c'est cette structure qu'il interroge ensuite pour elle-même.

Les noms modernes à chercher sont ceux des vingt-sept droites et des
double-six de Schläfli, du quadrangle généralisé d'ordre $(2,4)$, du groupe de
Weyl de $E_6$, des sous-systèmes de racines, des poids minuscules, et, pour la
fin, des carrés latins, des quasigroupes et de leurs autotopies.

\keywords{27 lines on a cubic surface, Schläfli graph, generalized quadrangle, double-six, Weyl group W(E6), E6 root system, del Pezzo surface, exceptional curves, Weyl group of D5, Weyl group of F4, quaternion group, quadratic form over F2, central extension, Baer sum, trihedral pair, Paley graph, affine plane of order 3, independent set, strongly regular graph, minuscule weight, highest root, Latin square, quasigroup, autotopism, Reidemeister condition, bitorsor, Freudenthal quartic}
\end{resume}

\pagerange{2}{119}
\section*{Le fil du dossier, et les conventions}

\subsection*{Les feuillets}

Le dossier compte cent dix-neuf pages. La page 1 est une couverture qui ne
porte qu'un « (34p) » au crayon ; la page 40 un intercalaire rose
(« (23 p) ») ; les pages 64, 65, 72, 76, 88, 93, 94, 99, 104 et 106 des
intercalaires vierges. La page 83 est dactylographiée : c'est l'annonce, par
l'Université des Sciences et Techniques du Languedoc, d'un enseignement de
DEUG A première année, option mathématique, pour 1976, intitulé
« Polyèdres réguliers », en deux parties confiées à A.~Grothendieck et à
C.~Contou-Carrère ; rien de sa main n'y figure, et sa date s'accorde avec celle
des archivistes. Les notes d'un cours de ce titre sont au dossier 86. Le
titre de l'inventaire annonce aussi une lettre ; les transcriptions n'en
trouvent pas. À partir de la page 95, les numéros que les archivistes ont
portés au crayon ont deux unités d'avance sur ceux du fac-similé, qu'on suit
ici comme le font les transcriptions.

Il pagine lui-même deux suites : les pages 2 à 39 (de « 1 » à « 32 », avec des
feuillets « bis » et « ter »), et les pages 41 à 57 (de « 1 » à « 17 »), qui
portent le titre « Les 27 droites de la surface cubique, et le syst.\ de
racines $E_6$ ». Ses renvois (« p.~8 et 9 », « th.\ p.~1 ») sont à ces
paginations ; on les traduit ici en pages du dossier.

\subsection*{Les stations}

\begin{itemize}
\item \textbf{Pages 2 à 10 : le théorème de caractérisation.} Un graphe à 27
  sommets, régulier de valence 10, où toute arête est dans un unique
  triangle, est le graphe cubique. Démonstration par un triangle
  $T = \lbrace u, v, w \rbrace$, l'involution $\sigma_T$, les quotients $P_u$,
  la structure de torseur des triangles extérieurs (lemme 7) et un
  étiquetage final (« Épilogue »).
\item \textbf{Pages 11 à 13 : dénombrements et un modèle.} La tour de
  sous-graphes du point au bitriangle, $|W| = 2^7 3^4 5 = 51\,840$, et une
  construction explicite du graphe à partir de deux ensembles à trois
  éléments.
\item \textbf{Pages 14 et 15 : l'algèbre des revêtements doubles.}
  Extensions centrales, produit de Baer sur $\Psi(V)$, formes quadratiques
  anisotropes ; existence si et seulement si $\dim V \leq 2$.
\item \textbf{Page 16 : les axiomes en valence quelconque.}
\item \textbf{Pages 17 à 39 : le graphe muni d'une configuration marquée.}
  Sommet (17 à 22), arête (23 à 25), triangle (25 à 31), balise et sablier
  (32, 33), prisme et bitriangle (34 à 39) : la donnée équivalente, le groupe
  de stabilité comme groupe de Weyl, les racines orthogonales.
\item \textbf{Pages 41 à 60 : les vingt-sept droites et $E_6$.} Le plan
  éclaté en six points, le groupe de Picard, la forme d'intersection, la
  classe $\lambda$ ; les racines, les réflexions, $W = \operatorname{Aut} C$,
  le théorème « $R(C)$ est de type $E_6$ » ; une définition intrinsèque et la
  liste des 72 bases.
\item \textbf{Pages 61 à 80 : un recensement.} Table des types de
  configurations ; parties libres ; bitriangles et leur auto-complémentarité ;
  polygones et polyèdres.
\item \textbf{Pages 81 à 98 : retour aux racines.} Différences et sommes de
  sommets, positions relatives de deux racines, les 27 en coordonnées dans
  les racines simples, et la recherche des ensembles $\mathfrak{S}$,
  $\mathfrak{S}_0$ qui devraient rendre les 27 à partir de $E_6$.
\item \textbf{Pages 100 à 118 : le triangle comme structure algébrique.}
  Trois sous-groupes distingués ; trois torseurs et une correspondance
  $\Gamma \subset P_1 \times P_2 \times P_3$ (un carré latin) ; ses
  autotopies ; les bitorseurs ; et, page 113, l'axiome qui revient aux pages
  7 et 8.
\item \textbf{Apartés.} Page 105 (le $56$ de $E_7$), page 119 (un
  autre sujet, $S^2 \times S^2$ et des idempotents).
\end{itemize}

\subsection*{Conventions, valables pour tout le dossier}

\begin{itemize}
\item Le graphe cubique est noté $C$ partout\footnote{Les pages l'appellent
  $S$ (pages 2 à 10 et 16), $C$ (à partir de la page 11) et $\Delta$ (pages 45
  à 60), et le disent tantôt « graphe », tantôt « complexe » cubique. On réserve
  $\Delta$ au \emph{bitriangle}, qui porte ce nom aux pages 34 à 39 et 73 à 75
  ; le graphe standard des pages 46 et 47, son $\Delta_0$, est ici $C_0$.}.
  Deux sommets sont \emph{liés} quand les droites correspondantes se
  rencontrent ; $E(x)$ est l'ensemble des sommets liés à $x$.
\item Une \emph{base} est un ensemble de six sommets deux à deux non liés
  (six droites deux à deux disjointes : l'une des deux moitiés d'un
  \emph{double-six} de Schläfli) ; à toute base $b$ correspond une unique
  \emph{base associée} (ou opposée) $b'$, et la paire $\lbrace b, b'\rbrace$ est
  une \emph{bibase}, c'est-à-dire un double-six.
\item $W = \operatorname{Aut}(C)$. On verra (page 52) qu'il est engendré par
  des réflexions et isomorphe au groupe de Weyl $W(E_6)$, d'ordre
  $51\,840 = 2^7 \cdot 3^4 \cdot 5$.
\item La forme bilinéaire est la forme d'intersection, avec ses signes : un
  sommet $\delta$ vérifie $\delta \cdot \delta = -1$, une racine $r \cdot r =
  -2$, et l'orthogonal $E$ de $\lambda$ est défini négatif. La classe
  $\lambda = 3\eta - \xi$ est l'opposée de la classe canonique.
\item Autour d'un triangle $T = \lbrace u, v, w\rbrace$ : $E(u)^* = E(u)
  \smallsetminus \lbrace v, w\rbrace$, de même pour $v, w$ ; $C^* = C
  \smallsetminus T$ ; $\sigma = \sigma_T$ l'involution de la page 2 ; $P_u =
  E(u)^*/\sigma$ ; $\widetilde\Gamma$ l'ensemble des 32 triangles contenus dans
  $C^*$ et $\Gamma = \widetilde\Gamma/\sigma$. L'extension de la page 14 est
  notée $\widetilde\Psi$\footnote{La page 14 la note $\widetilde\Gamma$ et
  appelle $\Gamma$ le groupe $\Psi(G)$ ; on garde $\Gamma$, $\widetilde\Gamma$
  pour les triangles, comme aux pages 7 et 8, et on suit la page 26, qui écrit
  $\widetilde\Psi(V, t)$.}.
\end{itemize}

Deux homonymies sont à garder en tête. Les « triades » de la page 61 sont des
partitions des 27 sommets en trois bitriangles ; celles de la page 73 sont des
parties de trois sommets deux à deux non liés d'un bitriangle. Et le symbole
$\Gamma$ des pages 107 à 116 est une correspondance entre trois torseurs ; il
n'est l'ensemble $\Gamma$ des triangles que dans le cas particulier de la
page 113, où les deux se rejoignent.

\pagerange{2}{10}
\section*{1. Le théorème de caractérisation (pages 2 à 10)}

\subsection*{L'énoncé}

\textbf{Théorème} (page 2). Soit $C$ un graphe (simple) tel que

\begin{enumerate}
\item[a)] $C$ a 27 sommets ;
\item[b)] tout sommet est lié à exactement 10 sommets ;
\item[c)] deux sommets liés ont exactement un voisin commun.
\end{enumerate}

Alors $C$ est isomorphe au graphe des 27 droites d'une surface cubique lisse
(le « graphe cubique »), et celui-ci vérifie a), b), c).

La condition c) dit que toute arête est dans un unique triangle. Le théorème
est vrai\footnote{On l'a vérifié dans le sens « seulement si » sur le modèle
standard des pages 45 et 46. Dans le langage d'aujourd'hui, les sommets et les
triangles de $C$ forment un \emph{quadrangle généralisé} d'ordre $(2, 4)$ : le
lemme 1 ci-dessous est exactement l'axiome des quadrangles (un point hors
d'une droite est collinéaire à un unique point de celle-ci), et le théorème
équivaut à l'unicité du quadrangle généralisé $\mathrm{GQ}(2, 4)$, ou du graphe
fortement régulier de paramètres $(27, 10, 1, 5)$, résultat classique qu'on
attribue en général à Seidel (1968). Ces noms sont de nous ; rien n'indique
qu'il ait connu cette littérature.}.

\subsection*{Le triangle et l'involution (pages 2 et 3)}

Soit $T = \lbrace u, v, w\rbrace$ un triangle (il en existe : une arête et son
troisième sommet).

\textbf{Lemme 1.} Tout sommet de $C^* = C \smallsetminus T$ est lié à un et un
seul sommet de $T$ ; $C^*$ est réunion disjointe de $E(u)^*$, $E(v)^*$,
$E(w)^*$, chacun de cardinal 8.

En effet $E(u)^* = E(u) \smallsetminus \lbrace v, w\rbrace$ a 8 éléments, et
$E(u)^* \cap E(v)^* = \emptyset$ puisque $E(u) \cap E(v) = \lbrace w\rbrace$
par c) ; les trois parties disjointes remplissent les $24 = 27 - 3$ sommets de
$C^*$.

\textbf{Lemme 2 et corollaire.} Pour $x \in E(u)^*$, il existe un unique $x'
\in E(u)^*$ lié à $x$ : le troisième sommet du triangle d'arête $\lbrace u,
x\rbrace$\footnote{La page écrit « $\exists!\, x' \in E(v)^*$ lié à $x$ », avec
$x'$ d'une lecture incertaine. Ainsi écrit, l'énoncé serait faux ($x$ a quatre
voisins dans $E(v)^*$, lemme 3) ; la démonstration (« reformulation de c) ») et
le corollaire qui suit portent sur le voisin de $x$ dans $E(u)^*$.}. D'où une
involution $\sigma = \sigma_T$ de $C$, qui fixe $u, v, w$ et envoie $x \in
E(u)^*$ sur ce sommet, de même sur $E(v)^*$ et $E(w)^*$ : chacun de ces trois
ensembles est partagé en quatre paires de sommets liés.

\textbf{Lemme 3.} Soit $x \in E(u)^*$. Alors $x$ est lié à exactement un
sommet de chaque paire $\lbrace y, \sigma y\rbrace$ de $E(w)^*$, donc à
exactement 4 sommets de $E(w)^*$ (et de même pour $E(v)^*$).

Il ne peut être lié aux deux, sinon l'arête $\lbrace y, \sigma y\rbrace$ serait
dans deux triangles, $\lbrace w, y, \sigma y\rbrace$ et $\lbrace x, y, \sigma
y\rbrace$ ; il est donc lié à au plus 4 sommets de $E(w)^*$ et à au plus 4 de
$E(v)^*$ ; or ses voisins hors de $\lbrace u, \sigma x\rbrace$ sont au nombre
de 8 et tous dans $E(v)^* \cup E(w)^*$.

\textbf{Corollaire.} Si $x \in E(u)^*$ et $y \in E(w)^*$ sont liés, $\sigma x$
et $\sigma y$ le sont : $\sigma_T$ est un automorphisme de $C$.

\subsection*{Les quotients $P_u$ et l'invariant $\delta_v$ (pages 3 à 7)}

Posons $P_u = E(u)^*/\sigma$, ensemble à quatre éléments, et de même $P_v$,
$P_w$. Pour $x \in E(u)^*$, la partie $\omega_v(x) = E(x) \cap E(v)^*$ choisit
un élément dans chacune des quatre paires : c'est une \emph{section} du
revêtement double $E(v)^* \to P_v$, et $\omega_v(\sigma x)$ est la section
complémentaire. Deux sections diffèrent par une fonction $\delta_v(x, y) \in
\mathbb{F}_2^{P_v}$, de support l'ensemble des paires où elles diffèrent ;
$\delta_v(x, \sigma y) = \delta_v(x, y) + \mathbf{1}$.

\textbf{Lemme 4.} Si $x, y \in E(u)^*$ sont distincts et non liés, ils ont
exactement deux voisins communs dans $E(v)^*$ ; autrement dit $\delta_v(x,
y)$ a un support de cardinal 2.

La démonstration (page 4) passe par le \textbf{lemme 5} : si $x \in E(u)^*$ et
$s \in E(w)^*$ sont liés, le troisième sommet $\xi$ du triangle $\lbrace x,
s\rbrace$ est dans $E(v)^*$ ; et si $x, y \in E(u)^*$ sont tous deux liés à
$s, t \in E(w)^*$, les troisièmes sommets $\xi$ de $(x, s)$ et $\xi'$ de $(y,
t)$ sont liés. S'il y avait trois voisins communs $s, t, r$, l'arête $\lbrace y,
\xi'\rbrace$ serait dans deux triangles.

Les paires de sommets non liés de $E(u)^*$ sont au nombre de $\binom 82 - 4 =
24$. Le lemme 4 fournit, pour chacune, la paire de ses deux voisins communs
dans $E(v)^*$ --- qui sont non liés, sinon on aurait deux triangles sur une
même arête ---, d'où une bijection
\[
\widetilde\varphi_{vu} : \mathfrak{P}'_2(E(u)^*) \longrightarrow
\mathfrak{P}'_2(E(v)^*)
\]
entre ensembles de paires non liées (« corollaire 2 », page 5), compatible à
$\sigma$. Passant aux quotients, elle induit une bijection canonique $c_{vu}$
de l'ensemble $\Pi_{22}(P_u)$ des trois partitions de $P_u$ en deux paires sur
$\Pi_{22}(P_v)$ ; le point à vérifier, que le résultat ne dépend que de la
partition et non de la paire choisie, est fait page 5.

\textbf{Lemme 6.} $c_{wv} \circ c_{vu} = c_{wu}$ ; plus précisément
$\widetilde\varphi_{wv} \circ \widetilde\varphi_{vu} = \sigma \circ
\widetilde\varphi_{wu}$ (pages 6 et 7). Les $c$ forment donc un système
transitif de bijections entre les trois ensembles $\Pi_{22}(P_u)$,
$\Pi_{22}(P_v)$, $\Pi_{22}(P_w)$.

\subsection*{Les triangles extérieurs forment un torseur (pages 7 et 8)}

On identifie entre eux les trois $\Pi_{22}$ grâce au lemme 6 : on obtient un
ensemble à trois éléments, que la page note d'un crochet $\sqsupset$. Un
ensemble à quatre éléments $P$ est canoniquement un torseur sous son groupe de
Klein (l'identité et les trois doubles transpositions), et les trois éléments
non nuls de ce groupe correspondent aux trois partitions $\Pi_{22}(P)$. Soit
donc $V$ l'espace vectoriel de dimension 2 sur $\mathbb{F}_2$ dont les
éléments non nuls sont $\sqsupset$ : $P_u$, $P_v$, $P_w$ deviennent trois
torseurs sous le \emph{même} $V$\footnote{La page écrit $\sqsupset = V^*$ et
« torseurs sous $V^*$ » : $V^*$ y désigne $V \smallsetminus \lbrace 0\rbrace$,
comme partout dans le dossier, et non le dual ; le groupe structural est
$V$.}.

Un triangle contenu dans $C^*$ a un sommet dans chacun des $E(\cdot)^*$ (un
sommet de $E(u)^*$ n'a qu'un voisin dans $E(u)^*$, son conjugué par $\sigma$, et
celui-ci n'est pas dans le même triangle que deux sommets de $E(u)^*$). Il est
déterminé par son sommet $x \in E(u)^*$ et son sommet dans
$\omega_v(x)$\footnote{La page écrit un $\prod$ ; le cardinal $32 = 8 \cdot 4$
qu'elle en tire est celui de la somme disjointe.} :
\[
\widetilde\Gamma \xrightarrow{\ \sim\ } \coprod_{x \in E(u)^*} \omega_v(x),
\qquad \operatorname{card} \widetilde\Gamma = 8 \cdot 4 = 32 .
\]
$\sigma$ opère sans point fixe sur $\widetilde\Gamma$ ; soit $\Gamma =
\widetilde\Gamma/\sigma$, de cardinal 16.

\textbf{Lemme 7.} L'application $\Gamma \to P_u \times P_v \times P_w$ est
injective, et son image est un torseur sous
\[
\Psi(V) = \lbrace (x, y, z) \in V^3 \mid x + y + z = 0\rbrace .
\]

L'injectivité est directe. Pour la structure de torseur, il suffit de voir
l'image stable par les sous-groupes $N_u = \operatorname{Ker}(\Psi \to V,
\mathrm{pr}_1)$, $N_v$, $N_w$, qui engendrent $\Psi$ ; elle aura alors $16 =
\operatorname{card} \Psi$ éléments. C'est le corollaire de la page 8 : si
$(x, s, \xi)$ est un triangle, $t \in E(v)^*$ un sommet non lié à $s$ mais lié
à $\xi$, et $\eta$ le second voisin commun de $s$ et $t$ dans $E(w)^*$, alors
$(x, \sigma t, \sigma \eta)$ est un triangle.

La page conclut par un programme : la donnée de $V$ et du torseur $\Gamma$
redonne $P_u$, $P_v$, $P_w$ et les triangles « à $\sigma$ près » ; il reste à
reconstituer les revêtements doubles $E(u)^* \to P_u$, \ldots,
$\widetilde\Gamma \to \Gamma$ et leurs liens. C'est ce que font l'algèbre des
pages 14 et 15 et, à la fin du dossier, la page 113.

\subsection*{Épilogue : l'étiquetage (pages 9 et 10)}

Fixons un triangle $(\alpha_4, \beta_4, \gamma_4)$ de $\widetilde\Gamma$, avec
$\alpha_4 \in E(u)^*$, $\beta_4 \in E(v)^*$, $\gamma_4 \in E(w)^*$, et posons
\[
\begin{gathered}
\lbrace \alpha_1, \alpha_2, \alpha_3, \alpha_4\rbrace = E(\beta_4) \cap
E(u)^*, \qquad \lbrace \beta_1, \ldots, \beta_4\rbrace = E(\gamma_4) \cap
E(v)^*, \\ \lbrace \gamma_1, \ldots, \gamma_4\rbrace = E(\alpha_4) \cap
E(w)^*,
\end{gathered}
\]
trois sections, et $\alpha'_i = \sigma\alpha_i$, etc. On a ainsi nommé les 24
sommets de $C^*$. Les tables de la page 9 en découlent :

\begin{itemize}
\item[(0)] $\alpha_i$ est lié à $\alpha'_i$, $\beta_i$ à $\beta'_i$, $\gamma_i$
  à $\gamma'_i$ ;
\item[(I)] les $\alpha_i$ sont liés à $\beta_4$, non à $\beta'_4$, et
  circulairement ; (I$'$) les $\alpha'_i$ sont liés à $\beta'_4$, non à
  $\beta_4$\footnote{Dans la première ligne de (I$'$), la page porte
  « $\beta_4$, $\beta_4$ » ; les accents de ces tables sont d'une lecture
  incertaine. On écrit ce que le lemme 3 impose.} ;
\item[(II)] (lemme 8) pour $1 \leq i \leq 3$, $\alpha_i$ est lié à
  $\gamma'_4$ et non à $\gamma_4$ --- sinon l'arête $\lbrace \beta_4,
  \gamma_4\rbrace$ serait dans les deux triangles de sommets $\alpha_i$ et
  $\alpha_4$ ---, et circulairement ; (II$'$) de même avec les primes
  échangés.
\end{itemize}

Ordonnons alors $\sqsupset$ ; par l'identification du lemme 6, cela numérote de
façon cohérente les trois éléments de $P_u \smallsetminus \lbrace
\pi_u(\alpha_4)\rbrace$, $P_v \smallsetminus \lbrace \pi_v(\beta_4)\rbrace$,
$P_w \smallsetminus \lbrace\pi_w(\gamma_4)\rbrace$. Avec cette numérotation
(lemme 9) :

\begin{itemize}
\item[(III)] pour $1 \leq i \neq j \leq 3$, $\alpha_i$ est lié à $\beta_j$,
  $\gamma_j$, $\beta'_i$, $\gamma'_i$, et non à $\beta'_j$, $\gamma'_j$,
  $\beta_i$, $\gamma_i$ ; et circulairement en $(\alpha, \beta, \gamma)$ ;
\item[(III$'$)] l'image de (III) par $\sigma$\footnote{La page n'écrit de
  (III$'$) que la première ligne, le reste en signes de répétition ; sa
  première lettre est lue $\alpha'_i$, peut-être $\alpha_i$. (III$'$) n'est de
  toute façon que l'image de (III) par l'automorphisme $\sigma$.}.
\end{itemize}

Le cœur de l'argument (page 10) : $\alpha_4$ est lié à $\beta'_1, \beta'_2,
\beta'_3, \beta_4$, et la section $\omega_v(\alpha_1)$ diffère de
$\omega_v(\alpha_4)$ sur exactement deux paires (lemme 4), dont celle de
$\beta_4$ exclue ; la partition correspondante est celle qui contient
$\lbrace 1, 4\rbrace$, d'où $\alpha_1$ lié à $\beta'_1$, $\beta_2$, $\beta_3$,
$\beta_4$. On a vérifié qu'une telle numérotation existe bien sur le modèle
standard.

\textbf{Conséquence.} Toutes les relations de liaison entre sommets de $C$
sont imposées par a), b), c) une fois faits des choix (triangle, triangle
extérieur, ordre sur $\sqsupset$) que tout graphe vérifiant ces conditions
permet de faire : deux tels graphes sont isomorphes, et le théorème est
démontré.

\pagerange{11}{12}
\section*{2. La tour des sous-graphes et l'ordre de $W$ (pages 11 et 12)}

Il considère une suite croissante de graphes $D_0 \subset D_1 \subset \cdots
\subset D_7$ : le vide, un point, une arête, un triangle, la
\emph{balise} (un triangle avec une arête pendante $uu'$), le
\emph{préprisme} (le triangle $u v w$ et une arête $u'v'$ avec $u'$ lié à $u$,
$v'$ à $v$), le \emph{prisme triangulaire}, et le \emph{bitriangle} (le
produit de deux triangles, neuf sommets, dix-huit arêtes : deux prismes
accolés par une face). Soit $\mathrm{Pl}_i$ l'ensemble des plongements
\emph{pleins} (comme sous-graphes induits) de $D_i$ dans $C$.

\begin{enumerate}
\item[a)] Les restrictions $\mathrm{Pl}_{i+1} \to \mathrm{Pl}_i$ sont
  surjectives, de fibres de cardinaux $27, 10, 1, 8, 4, 1, 1$ ;
\item[b)] d'où les cardinaux $27, 270, 270, 2160, 8640, 8640, 8640$ ;
\item[c)] chaque $\mathrm{Pl}_i$ est un espace homogène sous $W$.
\end{enumerate}

Les fibres se lisent sur le lemme 1 : pour la balise, $u' \in E(u)^*$ (8
choix) ; pour le préprisme, $v' \in E(v)^* \cap E(u')$ (4 choix, lemme 3) ; le
troisième sommet $w'$ du triangle $u'v'$ est forcé, et il est lié à $w$
puisque tout triangle de $C^*$ a un sommet dans chaque $E(\cdot)^*$ ; le
bitriangle se complète de même. En divisant par l'ordre du groupe
d'automorphismes de chaque $D_i$ on compte les configurations :
\[
\begin{array}{lllllll}
\text{points} & \text{arêtes} & \text{triangles} & \text{balises} &
\text{préprismes} & \text{prismes} & \text{bitriangles} \\
27 & 135 & 45 & 1080 & 4320 & 720 & 120
\end{array}
\]
avec les groupes $e$, $\mathfrak{S}_2$, $\mathfrak{S}_3$, $\mathfrak{S}_2$,
$\mathfrak{S}_2$, $\mathfrak{S}_2 \times \mathfrak{S}_3$ (ordre 12) et
$\mathfrak{S}_2 \cdot (\mathfrak{S}_3 \times \mathfrak{S}_3)$ (ordre
72)\footnote{Les 45 triangles sont les plans tritangents de la surface ; les
120 bitriangles sont les \emph{paires trièdres} de Steiner (deux triplets de
plans tritangents découpant les mêmes neuf droites). Noms de nous.}.

\textbf{d)} (page 12) Le stabilisateur dans $W$ d'un prisme plongé est
isomorphe à $\mathfrak{S}_3$ : il opère sur l'ensemble à trois éléments
$I = E(u) \cap E(v') \smallsetminus \lbrace u', v\rbrace$, et cet $I$ est
canoniquement isomorphe aux ensembles analogues construits sur les autres
arêtes du prisme. Comme c) donne $W$ transitif sur $\mathrm{Pl}_6$,
\[
\operatorname{card} W = 8640 \cdot 6 = 2^7 \cdot 3^4 \cdot 5 = 51\,840 .
\]
La page s'arrête sur un mot, « Compte »\footnote{Lecture incertaine, récrite
sur deux mots biffés, peut-être « Complexe cubique ».}. Les assertions c) et d)
ne sont pas démontrées sur ces feuillets ; elles le sont en substance par la
théorie des racines des pages 49 à 53 ($W$ simplement transitif sur les bases
ordonnées) et on les a vérifiées sur le modèle.

\pagerange{13}{13}
\section*{3. Un modèle explicite (page 13)}

La catégorie des graphes cubiques munis d'un triangle $I$ et d'un triangle
disjoint, pour les isomorphismes, équivaut à celle des couples $(I, \Pi)$ de
deux ensembles à trois éléments\footnote{La page écrit $\Pi'$ pour le second
triangle, lecture incertaine, et $\Pi$ pour l'ensemble abstrait.}. À un tel
couple on associe, avec $\omega(I)$ l'ensemble (à deux éléments) des ordres
circulaires sur $I$, chaque $\rho \in \omega(I)$ vu comme une permutation
circulaire,
\[
C = \bigl(I \times \lbrace 0, 1, 2\rbrace\bigr) \amalg \bigl(I \times \Pi
\times \omega(I)\bigr) \qquad (9 + 18 = 27 \text{ sommets}),
\]
avec les liaisons suivantes (entre sommets distincts) :

\begin{itemize}
\item $(i, \nu)$ et $(j, \mu)$ sont liés si et seulement si $i = j$ ou $\nu =
  \mu$ (un bitriangle) ;
\item $(i, \nu)$ et $(j, \pi, \rho)$ sont liés si et seulement si : $i = j$ et
  $\nu = 0$, ou $i \neq j$, $\nu = 1$ et $j = \rho(i)$, ou $i \neq j$, $\nu =
  2$ et $i = \rho(j)$ ;
\item $(j, \pi, \rho)$ et $(j', \pi', \rho')$ sont liés si et seulement si :
  pour $j = j'$, $\pi = \pi'$ (et donc $\rho \neq \rho'$) ; pour $j \neq j'$,
  ou bien $\pi \neq \pi'$ et $\rho = \rho'$, ou bien $\pi = \pi'$ et $\rho \neq
  \rho'$\footnote{Dans le dernier cas la page écrit « $i \neq j$, $\rho =
  \rho'$ ; $i = j$, $\rho \neq \rho'$ », alors que seuls $j, j'$ figurent dans
  le couple ; on lit $\pi \neq \pi'$, $\pi = \pi'$. On a vérifié que, avec
  cette lecture, le graphe obtenu est fortement régulier de paramètres $(27,
  10, 1, 5)$, donc le graphe cubique ; l'autre lecture naturelle ne l'est
  pas.}.
\end{itemize}

L'automorphisme $\sigma$ du triangle $I \times \lbrace 0\rbrace$ est
l'identité sur $I$, échange $1$ et $2$ sur le premier morceau, et échange les
deux ordres circulaires sur le second ; on vérifie que c'est bien
l'involution $\sigma_T$ de la page 2 pour $T = I \times \lbrace 0\rbrace$. Le bas
de la page esquisse d'autres écritures du même ensemble, avec des
$\mathbb{F}_2$-torseurs ; elles ne sont pas menées.

\pagerange{14}{15}
\section*{4. L'algèbre des revêtements doubles (pages 14 et 15)}

La structure qui manquait à la fin du lemme 7 --- les revêtements doubles
$\widetilde\Gamma \to \Gamma$, $E(u)^* \to P_u$ --- relève d'un énoncé
d'algèbre que voici, dans le cas qui sert, celui d'un $\mathbb{F}_2$-espace
vectoriel\footnote{La page énonce a) et b) pour un groupe abélien $G$
quelconque, avec $V = G \otimes \mathbb{F}_2$. Dans ce cadre, $p_j \alpha_i$
vaut $\mathrm{id}_G$ pour un indice $j \neq i$ et $-\mathrm{id}_G$ pour
l'autre, et la classe d'une extension centrale de $G$ par $\varepsilon$ n'est
pas en général une application quadratique ; on n'a pas vérifié la version
générale et on s'en tient à $2G = 0$, qui est l'hypothèse de c) et le seul cas
utilisé.}.

Soit $V$ un $\mathbb{F}_2$-espace vectoriel de dimension finie, $\varepsilon$
un $\mathbb{F}_2$-espace vectoriel, $\Psi = \Psi(V) = \operatorname{Ker}(V^3
\xrightarrow{\ +\ } V)$, $p_i : \Psi \to V$ les projections et $\alpha_i : V
\to \Psi$ les injections d'image $N_i = \operatorname{Ker} p_i$ (par exemple
$\alpha_1(g) = (0, g, g)$), de sorte que $p_j \alpha_i$ est nul si $i = j$ et
l'identité sinon.

\textbf{a)} Soit $1 \to \varepsilon \to E_0 \to V \to 1$ une extension
centrale ; sa classe est la forme quadratique $q : V \to \varepsilon$, $q(g) =
\tilde g^2$, et toute forme quadratique provient d'une telle
extension\footnote{C'est l'isomorphisme $H^2(V, \varepsilon) \simeq
\operatorname{Quad}(V, \varepsilon)$ pour $V$ abélien élémentaire, la page
écrit $\operatorname{Quad}_{\mathbb{Z}}(G, \varepsilon) \simeq
\operatorname{Quad}_{\mathbb{F}_2}(V, \varepsilon)$.}. Soit
\[
\widetilde\Psi = p_1^* E_0 \wedge p_2^* E_0 \wedge p_3^* E_0
\]
le produit de Baer des trois images inverses, extension centrale de $\Psi$ par
$\varepsilon$\footnote{« produit de Baer » : lecture incertaine du second
mot, confirmée par le « produit de Baer » de la page 37.}. Pour chaque $i$,
$\alpha_i^* \widetilde\Psi \simeq E_0 \wedge E_0$ est canoniquement scindée
($g \mapsto [\tilde g, \tilde g]$, indépendant du relèvement puisque
$\varepsilon$ est d'exposant 2), d'où des relèvements $\rho_i : V \to
\widetilde\Psi$ de $\alpha_i$, et
\[
\rho_1(g)\, \rho_2(g)\, \rho_3(g) = q(g) \qquad (g \in V),
\]
ce qu'on vérifie en écrivant $\rho_1(g) = [1, \tilde g, \tilde g]$, etc.

\textbf{b)} Inversement, tout système $(\widetilde\Psi, \rho_1, \rho_2,
\rho_3)$ d'une extension centrale de $\Psi$ par $\varepsilon$ et de
relèvements des $\alpha_i$ provient d'une extension $E_0$, déterminée à
isomorphisme unique près, de classe donnée par la formule précédente ; et la
catégorie de ces systèmes est rigide. On rapporte cet énoncé sans l'avoir
vérifié en entier.

\textbf{c)} Conditions équivalentes :
\begin{enumerate}
\item[(i)] $q$ est \emph{anisotrope} : $q(x) \neq 0$ pour $x \neq 0$ ;
\item[(ii)] posant $\widetilde P_i = \widetilde\Psi/\rho_i(V)$, trois points
  de $\widetilde\Psi$ qui ont deux à deux la même image dans le
  $\widetilde P_k$ convenable sont égaux\footnote{Les indices des égalités de
  (ii) sont d'une lecture incertaine ; on donne le sens qui rend (ii)
  équivalent à (iii).} ;
\item[(iii)] si $z_i \in \rho_i(V)$ et $z_1 z_2 z_3 = 1$, alors $z_1 = z_2 =
  z_3 = 1$ ;
\item[(iv)] si $\rho_1(g) \rho_2(g) \rho_3(g) = 1$, alors $g = 0$.
\end{enumerate}
(iii) et (iv) sont équivalents parce que $\rho_1(a) \rho_2(b) \rho_3(c)$ est
au-dessus de $(b + c, a + c, a + b)$, nul seulement si $a = b = c$ ; et (iv)
est (i) par la formule de a).

\textbf{Existence.} Pour $\varepsilon = \mathbb{F}_2$\footnote{La page porte
« $\varepsilon$ $\neq$ $\mathbb{F}_2$ », le signe étant d'une lecture
incertaine ; seul $\varepsilon = \mathbb{F}_2$ donne l'énoncé qui suit.}, un
tel système existe si et seulement si $\dim V \leq 2$, et il est alors unique à
isomorphisme unique près. C'est juste : une forme quadratique sur
$\mathbb{F}_2^n$ à valeurs dans $\mathbb{F}_2$ est isotrope dès que $n \geq 3$
(théorème de Chevalley--Warning), et en dimension 2 la seule forme anisotrope
est $x^2 + xy + y^2$. L'extension $E_0$ correspondante de $V = \mathbb{F}_2^2$
est alors le \emph{groupe des quaternions} $Q_8$ : tout élément non nul de $V$
se relève en un élément d'ordre 4\footnote{Le nom est de nous ; la page 26
décrit la même extension comme « l'extension de $V$ par $\mathbb{F}_2$ qui ne
splitte sur aucun sous-groupe $\simeq \mathbb{F}_2$ de $V$ ».}.

Pour $\dim V = 2$, $\widetilde\Psi$ a $2 \cdot 16 = 32$ éléments et chaque
$\widetilde P_i$ en a 8 : ce sont les cardinaux de $\widetilde\Gamma$ et des
$E(u)^*$. C'est ainsi que les pages 25 et 26 reconstruiront $C$ autour d'un
triangle.

\pagerange{16}{16}
\section*{5. Les axiomes en valence quelconque (page 16)}

Soit $S$ un graphe non vide tel que a) deux sommets liés ont un unique voisin
commun, b) $S$ est régulier de valence $\nu \geq 1$.

\textbf{Lemme 1.} $\nu$ est pair, et $\operatorname{card} S \geq 3(\nu - 1)$,
avec égalité si et seulement si, pour un (ou tout) triangle $T$, tout sommet
hors de $T$ est lié à exactement un sommet de $T$.

Les voisins d'un sommet se groupent en paires (les triangles qui le
contiennent), d'où la parité ; les $\nu - 2$ voisins extérieurs de chaque
sommet d'un triangle $T$ sont deux à deux distincts d'un sommet à l'autre, d'où
$\operatorname{card} S \geq 3 + 3(\nu - 2)$\footnote{L'inégalité est récrite
sur un autre chiffre ; c'est bien $3(\nu - 1)$ qu'il faut, et pour $\nu = 10$
on retrouve les 27 sommets de la page 2.}. Les calculs griffonnés qui suivent
posent $\nu = 2(\mu + 1)$, $\operatorname{card} S = 3(2\mu + 1)$ dans le cas
d'égalité ; les non-voisins d'un sommet sont alors $4\mu$ (16 pour $\mu = 4$).
L'inégalité $4\mu \leq 2^{\mu + 1}$ qui suit s'explique si l'on code un
non-voisin $x$ de $s$ par la section qu'il choisit dans les $\mu + 1$ paires de
$E(s)$, comme le fait la page 17 ; c'est notre lecture, la page ne le dit pas.

\pagerange{17}{20}
\section*{6. Un sommet marqué (pages 17 à 20)}

Soit $u \in C$ et $Q_u$ l'ensemble à cinq éléments des triangles de sommet $u$
(les paires de sommets liés de $E(u)$)\footnote{La lettre de cet ensemble
tient, page 17, du $P$ et du $Q$ cursifs ; la page 18 l'écrit nettement
$Q_u$.}. Soit $\Psi = \Psi(Q_u) = \operatorname{Ker}(\mathbb{F}_2^{Q_u}
\xrightarrow{\ \Sigma\ } \mathbb{F}_2)$, de cardinal 16.

\textbf{Description.} La donnée de $(C, u)$ équivaut à celle d'un ensemble
$Q_u$ à cinq éléments et d'un torseur $P$ sous $\Psi(Q_u)$. Dans un sens :
$E(u)$, muni de son involution sans point fixe, est un revêtement double de
$Q_u$, c'est-à-dire un torseur $Q$ sous $\mathbb{F}_2^{Q_u}$ (l'ensemble de ses
32 sections) ; tout $x \notin E(u) \cup \lbrace u\rbrace$ choisit un élément
dans chaque paire (lemme 1), donc une section ; l'application $C
\smallsetminus (E(u) \cup \lbrace u\rbrace) \to Q$ est injective, d'image un
torseur $P$ sous $\Psi$ (16 éléments). Dans l'autre :
\[
C \simeq \lbrace u\rbrace \amalg \coprod_{i \in Q_u} \bigl(P \wedge^{\Psi}
(\mathbb{F}_2, p_i)\bigr) \amalg P ,
\]
avec les règles a) $u$ est lié aux $10$ points du terme médian, non à ceux de
$P$ ; b) les deux points de chaque $P \wedge^\Psi (\mathbb{F}_2, p_i)$ sont
liés, et non liés aux autres points du terme médian ; c) $a \in P \wedge^\Psi
(\mathbb{F}_2, p_i)$ est lié à $x \in P$ si et seulement s'il est l'image de
$x$ ; d) $x, y \in P$ sont liés si et seulement si $y - x$ est l'un des cinq
vecteurs $e'_i = \sum_{j \neq i} e_j$ de poids 4.

\textbf{Lemme} (page 18). Pour $x, y \notin E(u) \cup \lbrace u\rbrace$
distincts, $x$ et $y$ sont liés si et seulement s'il existe un et un seul
triangle $t$ de sommet $u$ sur lequel $x$ et $y$ font le même
choix\footnote{La page écrit « $t - E(x) = t - E(y)$ », et précise que les
deux membres sont réduits à un élément $\neq u$ ; il faut lire $t \cap E(x) =
t \cap E(y)$ (ou ôter $u$). On a vérifié la règle sur le modèle.}. Les
différences de sections sont de poids pair (elles sont dans $\Psi$) ; aucun
accord est impossible ($\sum e_i \notin \Psi$, « car $5 \neq 0$ dans
$\mathbb{F}_2$ ») ; un seul accord correspond au poids 4, trois accords au
poids 2.

\textbf{Corollaire} (page 19). Le stabilisateur $W_u$ de $u$ est le groupe des
couples $(\varphi, \bar\varphi)$ d'une permutation de $Q_u$ et d'un
automorphisme affine de $P$ au-dessus d'elle :
\[
1 \to \Psi(Q_u) \simeq \mathbb{F}_2^4 \to W_u \to \mathfrak{S}_5 \to 1 ,
\]
extension scindée. C'est le groupe de Weyl $W(D_5)$, d'ordre $16 \cdot 120 =
1920$, d'indice 27 dans $W \simeq W(E_6)$\footnote{Il écrit « $W_{D_5}$ de
$D_5$ ($\simeq SO(10)$) » : le groupe de Weyl du type $D_5$, celui de
$SO(10)$.}.

\subsection*{Les bases (pages 19 et 20)}

Il annonce un système de racines $D_5$ attaché à $(C, u)$ et le relie aux
« bases ». Il y a $\binom{10}{2} - 5 = 40$ paires de sommets non liés dans
$E(u)$.

\textbf{Lemme.} Pour toute paire $\lbrace \xi_1, \xi_2\rbrace$ de sommets non
liés de $E(u)$, il existe une unique base $b$ telle que $b \cap E(u) =
\lbrace\xi_1, \xi_2\rbrace$, et alors $b' \cap E(u) = \lbrace \sigma\xi_1,
\sigma\xi_2\rbrace$, où $\sigma$ est l'involution de $E(u)$.

L'existence vient de la transitivité de $W_u$ sur ces paires ; la preuve de
l'unicité est dans un bloc biffé, presque illisible, et n'est pas
reconstituée ici\footnote{Le bloc est encadré et barré de deux traits en
croix ; seuls quelques signes s'y lisent. L'énoncé est vrai (on l'a vérifié :
les 40 bases « orthogonales » à $u$ correspondent une à une aux 40 paires).}.

\textbf{Positions d'un sommet et d'une base} (page 20). Pour $u$ fixé, les 72
bases se répartissent en $16 + 16 + 40$ :
\begin{itemize}
\item[(i)] $u \in b$ (16 bases) : $E(u) \cap b = \emptyset$, $\operatorname{card}
  E(u) \cap b' = 5$, et $r_b \cdot u = 1$ ;
\item[(ii)] $u \in b'$ (16 bases) : les rôles de $b$ et $b'$ échangés, $r_b
  \cdot u = -1$ ;
\item[(iii)] $u \notin b \cup b'$ (40 bases, dites \emph{orthogonales} à $u$)
  : $\operatorname{card} E(u) \cap b = \operatorname{card} E(u) \cap b' = 2$,
  $r_b \cdot u = 0$.
\end{itemize}
Ici $r_b$ est la racine attachée à la bibase (page 49).

\pagerange{21}{22}
\section*{7. Le réseau d'un revêtement double, et $D_5$ (pages 21 et 22)}

Soit $\widetilde Q \to Q$ un revêtement double d'un ensemble à $r$ éléments,
d'involution $\sigma$. Sur $\mathbb{Z}^{\widetilde Q}$, posons $\tilde q \cdot
\tilde q = -1$, $\tilde q \cdot \sigma\tilde q = 1$, et $\tilde q \cdot \tilde
q' = 0$ si $\tilde q'$ n'est ni $\tilde q$ ni $\sigma \tilde q$\footnote{La page
écrit le produit sur $\mathbb{Z}^Q$, sans tilde.}. Soit $E$ le quotient par le
noyau de la forme. Comme $\tilde q + \sigma \tilde q$ est dans le noyau, $E$
est libre de rang $r$, de base un relèvement de chaque point de $Q$, la forme y
étant $-\mathrm{id}$ ; les « racines » $\tilde q - \tilde q'$ ($\tilde q' \neq
\tilde q, \sigma \tilde q$) sont les $\pm e_i \pm e_j$, un système de type
$D_r$.

Pour $Q = E(u)$ : dans une base où $u = \xi'_6$, $E(u)$ est formé des
$\xi_i$ et des $\xi_{i6}$ ($1 \leq i \leq 5$), $\sigma \xi_i = \xi_{i6}$, et

\begin{enumerate}
\item[a)] les $x + \sigma x$, $x \in E(u)$, sont tous égaux à $\theta = \eta -
  \xi_6 = \lambda - u$\footnote{La page écrit « $\theta = \eta - \xi'_6$ » (on
  lit aussi « $y - \xi_6$ »), et prend $\theta_0 = \tfrac12(\xi'_6 - \eta)$
  dans la définition des $\varepsilon_i$. Le calcul $\xi_i + \xi_{i6} = \eta -
  \xi_6$ impose $\theta = \eta - \xi_6$. L'écriture $\theta = \lambda - u$ est
  de nous : c'est la relation $x + \sigma x + u = \lambda$, la somme des trois
  droites d'un plan tritangent étant la section hyperplane (page 42).} ;
\item[b)] $x \mapsto x - \tfrac12 \theta$ envoie $E(u)$ sur les $\pm
  \varepsilon_i$, $\varepsilon_i = \xi_i - \tfrac12\theta$, qui vérifient
  $\varepsilon_i \cdot \varepsilon_j = -\delta_{ij}$ ; c'est un isomorphisme du
  réseau de la page 21 sur un sous-réseau de $\tfrac12 \widehat E(C)$, et les
  différences de sommets non liés de $E(u)$, inchangées par la translation,
  sont les 40 racines $\pm\varepsilon_i \pm \varepsilon_j$, de type $D_5$ :
  ce sont les racines de $E_6$ orthogonales à $u$.
\end{enumerate}

\pagerange{23}{25}
\section*{8. Une arête marquée (pages 23 à 25)}

Page 23, deux vérifications de l'ordre de $W$. Une base orthogonale à $u$
équivaut à un ensemble $I$ à six éléments ($b$ lui-même) muni d'une partie
$I_2 = b \cap E(u)$ à deux éléments ; d'où\footnote{La page écrit le membre
de droite « $27 \cdot 3^3 \cdot 5$ » (premier facteur lu aussi $2^7$) ; le
produit vaut $2^7 \cdot 3^4 \cdot 5$.}
\[
\operatorname{card} W = (27 \cdot 40) \cdot (2 \cdot 4!) = 51\,840 .
\]

Marquer une arête $\lbrace v, w\rbrace$ revient à marquer le triangle $\lbrace
u, v, w\rbrace$ avec son sommet $u$, donc un élément $t$ de $Q_u$ ; $P_u = Q_u
\smallsetminus \lbrace t\rbrace$ a 4 éléments et $\Psi(Q_u) \simeq
\mathbb{F}_2^{P_u}$ par projection. Un torseur sous
$\mathbb{F}_2^{P_u}$ est un ensemble $\widetilde P_u$ à huit éléments muni d'une
involution sans point fixe, et
\[
C \simeq \lbrace u\rbrace \amalg \lbrace v, w\rbrace \amalg \widetilde P_u
\amalg P , \qquad 1 + 2 + 8 + 16 = 27,
\]
avec six règles de liaison a)--f) analogues à celles de la page 18, la
dernière disant que $x, y \in P$ sont liés si et seulement si $x - y$ est l'un
des vecteurs $e'_i = \sum_{j \neq i} e_j$ ($i \in P_u$) ou $e'_0 = \sum_j e_j$.

\textbf{Corollaire.} Le stabilisateur $W_\alpha$ de l'arête $\alpha$ est une
extension scindée de $\mathfrak{S}_{P_u} \simeq \mathfrak{S}_4$ par
$\mathbb{F}_2^4$ : c'est $W(B_4)$, d'ordre 384, dont le système de racines a
32 racines de deux longueurs. Il le cherche dans $E_6$ ; ce qu'il trouve
(page 25), c'est qu'en \emph{épinglant} le triangle (en fixant ses trois
sommets) on restreint le groupe structural à $\Psi_{P_u} = \operatorname{Ker}(
\mathbb{F}_2^{P_u} \to \mathbb{F}_2)$, d'où
\[
1 \to \Psi_{P_u} \simeq \mathbb{F}_2^3 \to W_t' \to \mathfrak{S}_{P_u} \to 1,
\qquad W_t' \simeq W(D_4),
\]
d'ordre 192, où $W_t'$ est le groupe des automorphismes qui fixent $t$ point
par point ; le système $D_4$ a $\binom 82 - 4 = 24$ racines, toutes de même
longueur.

\pagerange{25}{31}
\section*{9. Un triangle marqué (pages 25 à 31)}

\textbf{Description} (pages 25 et 26). La donnée de $(C, t)$, $t$ un triangle,
équivaut à celle de

\begin{enumerate}
\item[a)] un ensemble $t$ à trois éléments ;
\item[b)] un $\mathbb{F}_2$-espace vectoriel $V$ de dimension 2 (ou, ce qui
  revient au même, l'ensemble $V^* = V \smallsetminus \lbrace 0\rbrace$ à trois
  éléments)\footnote{La lettre ressemble d'abord à un $N$ ; la page 26 écrit
  $V$.} ;
\item[c)] un torseur $\widetilde P$ sous l'extension canonique
  $\widetilde\Psi(V, t) = \bigwedge_{i \in t} p_i^*(\widetilde V)$ de
  $\Psi(V, t) = \operatorname{Ker}(V^t \to V)$, où $\widetilde V$ est
  l'extension de $V$ par $\mathbb{F}_2$ qui ne se scinde sur aucun sous-groupe
  d'ordre 2 --- le groupe $Q_8$ de la page 15.
\end{enumerate}

On pose $\widetilde P_i = \widetilde P/N'_i$, où $N'_i \subset \widetilde\Psi$
est le relèvement canonique de $N_i = \operatorname{Ker} p_i$ ; alors $C \simeq
t \amalg \coprod_{i \in t} \widetilde P_i$, avec a) les sommets de $t$ deux à
deux liés ; b) $i$ lié aux points de $\widetilde P_i$ et non aux autres ; c)
dans $\widetilde P_i$, $x$ et $y$ liés si et seulement si $y = \sigma x$ ; d)
pour $i \neq j$, $x \in \widetilde P_i$ et $y \in \widetilde P_j$ liés si et
seulement s'ils sont images d'un même point de $\widetilde P$. Ici $\widetilde
P$ est l'ensemble $\widetilde\Gamma$ des 32 triangles extérieurs à $t$, et les
$\widetilde P_i$ sont les $E(i)^*$ : c'est l'achèvement du programme de la
page 8.

\textbf{Corollaire.} Le groupe $W_t$ des automorphismes de $(C, t)$ est
l'ensemble des triples $(\tau, \rho, \varphi)$ d'une permutation de $t$, d'un
automorphisme de $V$ et d'un automorphisme de $\widetilde P$ au-dessus de
l'automorphisme induit de $\widetilde\Psi$ :
\[
1 \to \widetilde\Psi \to W_t \to \mathfrak{S}_t \times \mathfrak{S}_{V^*} \to 1,
\qquad \operatorname{card} W_t = 32 \cdot 36 = 1152 = 51\,840/45 .
\]
« Sauf erreur », dit-il, c'est le groupe de Weyl de $D_4$ complété par ses
automorphismes extérieurs, $\mathfrak{S}_t$ correspondant à ceux-ci, et la
catégorie des systèmes de racines de type $D_4$ serait équivalente à celle des
graphes cubiques munis d'un triangle\footnote{Le groupe $W(D_4) \rtimes
\mathfrak{S}_3$ est isomorphe au groupe de Weyl $W(F_4)$, d'ordre 1152, et
c'est le groupe des automorphismes du système de racines $D_4$ ; les ordres
s'accordent. L'équivalence de catégories est son énoncé, non vérifié ici. Le
nom $W(F_4)$ est de nous.}.

\textbf{Les 24 racines orthogonales à un triangle} (pages 27 et 28). Les bases
$b$ telles que $b \cap E(i)$ ait deux éléments pour chaque $i \in t$ --- ce
qui force $b \cap t = \emptyset$ --- sont au nombre de $8 \cdot 6/2 = 24$, et
$b \mapsto b \cap E(i)$ les met en bijection avec les paires non liées de
$\widetilde P_i$. Les racines correspondantes sont les racines orthogonales à
$t$ ; elles forment un système de type $D_4$ (« à vérifier », écrit-il ; c'est
exact). En marge : « $D_4$ = carrousel », avec un dessin : le triangle et deux
arêtes pendantes à chaque sommet.

Le système $(C, t, b)$, $b$ orthogonale à $t$, équivaut à une base $b$ munie
d'une partition en trois paires $\lbrace \xi_1, \xi_2\rbrace$, $\lbrace \xi_3,
\xi_4\rbrace$, $\lbrace \xi_5, \xi_6\rbrace$, le triangle étant $\lbrace
\xi_{12}, \xi_{34}, \xi_{56}\rbrace$ ; il y en a $15$ par base, soit $72 \cdot
15 = 1080 = 45 \cdot 24$, et le stabilisateur est $W(B_3) = \mathfrak{S}_2 \wr
\mathfrak{S}_3$, d'ordre 48 ; $1080 \cdot 48 = 51\,840$.

\textbf{Position d'un triangle et d'une base} (pages 29 à 31). Comme $u + v + w
= \lambda$ est orthogonal à toutes les racines, $\sum_{i \in t} i \cdot r_b =
0$ avec chaque terme dans $\lbrace -1, 0, 1\rbrace$ : ou bien $b$ est
orthogonale à $t$ (24 bases), ou bien, à l'ordre près, $r_b \cdot u = 1$ ($u
\in b$), $r_b \cdot v = -1$ ($v \in b'$), $r_b \cdot w = 0$ (48 bases, en six
groupes de huit selon les rôles de $u, v, w$). Dans le second cas, la donnée
de $(t, b)$ équivaut à celle d'un triangle épinglé et d'un point $u' \in
E(w)^*$ ($27 \cdot 10 \cdot 8 = 2160$), ou d'une base avec deux éléments
ordonnés $\xi_1, \xi_2$, le triangle étant $(\xi_1, \xi'_2, \xi_{12})$ ($72
\cdot 6 \cdot 5 = 2160$). Le stabilisateur est $\mathfrak{S}_4 = W(A_3) \simeq
W(D_3)$, « isomorphisme exceptionnel » $D_3 \simeq A_3$ ; ses 12 racines sont
les $\alpha - \beta$, $\alpha, \beta$ distincts non liés dans $E(w)^*
\smallsetminus \lbrace u', \sigma u'\rbrace$.

\textbf{Deux triangles disjoints et une base orthogonale aux deux} (page 31).
Ils correspondent à un ensemble $I$ à six éléments muni de deux partitions de
type $(2,2,2)$ sans classe commune, c'est-à-dire d'une structure
\emph{hexagonale} sur $I$ (les deux partitions sont les deux familles de côtés
alternés). Le triangle $t''$ qui complète $\lbrace t, t'\rbrace$ en bitriangle
est encore orthogonal à $b$. Pour $b$ fixée il y a $120 = 5!$ couples $(t,
t')$ et $60 = 5!/2$ paires, autant que de structures hexagonales sur un
ensemble à six éléments ; d'où $72 \cdot 5! = 45 \cdot 32 \cdot 6 = 8640 =
\operatorname{card} W/6$ triples, $6 = \operatorname{card}\mathfrak{S}_t$.

\pagerange{32}{33}
\section*{10. Balise et sablier (pages 32 et 33)}

Une \emph{balise} est un triangle $\lbrace u, v, w\rbrace$ avec une arête
pendante $\lbrace u, x\rbrace$. Sa donnée revient à celle d'un triangle à
sommet marqué (un torseur sous $\mathbb{F}_2^{P_u}$, page 23) et d'un point de
$\widetilde P_u$ ; il reste un ensemble $R_u$ à trois éléments et un
revêtement double de $R_u$. La catégorie des $(C, \text{balise})$ est discrète
au sens où ses objets n'ont pas d'automorphisme non trivial une fois la balise
épinglée ; le groupe d'automorphismes de $(C, \text{balise})$ est le groupe de
Weyl élargi de $D_3$, c'est-à-dire $W(B_3)$, d'ordre $3! \cdot 2^3 = 48$, et
\[
45 \cdot 3 \cdot 8 = 1080 = 51\,840 / 48 .
\]
Les racines correspondantes sont les $\alpha - \beta$, $\alpha, \beta \in
E(u)^* \smallsetminus \lbrace x, \sigma x\rbrace$ distincts non liés.

Un \emph{sablier} est une paire de triangles de sommet commun : la donnée
équivaut à un sommet $u$ et une partition de $Q_u$ en une partie à deux
éléments et une à trois. Il y en a $27 \cdot \binom 52 = 270$, de groupe
$(\mathfrak{S}_2 \times \mathfrak{S}_3) \cdot \Psi(Q_u)$ d'ordre $2 \cdot 6
\cdot 16 = 192$, et $270 \cdot 192 = 51\,840$. Un sablier épinglé équivaut à
une balise épinglée ($270 \cdot 8 = 1080 \cdot 2 = 2160$)\footnote{La page
dit que l'épinglage du sablier « élimine le groupe d'autom.\ des diagrammes,
qui est $\mathbb{Z}/4\mathbb{Z}$ ». Le groupe des automorphismes du graphe
formé de deux triangles à sommet commun est d'ordre 8 (diédral), et c'est
l'ordre 8 que demande l'égalité des deux comptes.}.

\pagerange{34}{39}
\section*{11. Prisme et bitriangle (pages 34 à 39)}

Le prisme triangulaire, dit-il, équivaut à un carré à côté marqué, et il passe
au \emph{bitriangle} $\Delta$, dont il y a 120 (page 11). La donnée d'un
bitriangle abstrait équivaut à celle d'un ensemble $\Delta_0$ à neuf éléments
muni de deux partitions de type $(3,3,3)$ telles que $\Delta_0 \to
\Delta_0/\Pi \times \Delta_0/\Pi'$ soit bijective ; ou à celle d'un ensemble
$T$ à six éléments (les six triangles de $\Delta$) muni d'une partition de type
$(3, 3)$, c'est-à-dire d'une application $T \to I$ de degré 3 sur un ensemble
$I$ à deux éléments, avec $\Delta = \prod_{i \in I} T_i$. D'où
$\operatorname{Aut}(\Delta) \simeq \mathfrak{S}_2 \ltimes (\mathfrak{S}_3
\times \mathfrak{S}_3)$, d'ordre 72.

Pour $\Delta$ plongé dans $C$, le noyau de $\operatorname{Aut}(C, \Delta) \to
\operatorname{Aut}(\Delta)$ est d'ordre 6, isomorphe à $\mathfrak{S}_3$ (le d)
de la page 12). Il l'interprète par un ensemble $V^*$ à trois éléments,
canoniquement isomorphe, pour chaque $u \in \Delta$, à l'ensemble $V^*(u)$ des
trois triangles de sommet $u$ non contenus dans $\Delta$, les isomorphismes
formant un système transitif ; d'où
\[
\operatorname{Aut}(C, \Delta) \xrightarrow{\ \sim\ } \operatorname{Aut}(\Delta)
\times \operatorname{Aut}(V^*), \qquad 432 = 72 \cdot 6 = 51\,840/120 .
\]

\textbf{Lemme 2} (page 36). a) Pour $x \in C \smallsetminus \Delta$,
$\Delta(x) = \Delta \cap E(x)$ a trois éléments et c'est le graphe d'une
bijection $T_i \simeq T_{i'}$ (une transversale du carré $3 \times 3$) ; b)
soit $\mathbb{F}'$ l'ensemble (à six éléments) de ces transversales\footnote{La
lettre est un F à double barre, rendue $\mathbb{F}'$ comme dans la
transcription.} ; l'application $C \smallsetminus \Delta \to \mathbb{F}' \times
V^*$ est bijective ($18 = 6 \cdot 3$), et
\[
C \simeq \Delta \amalg (\mathbb{F}' \times V^*) ,
\]
deux points distincts $x, y$ de $\mathbb{F}' \times V^*$ étant liés si et
seulement si ou bien $\Delta(x) \cap \Delta(y) \neq \emptyset$ et $v(x) = v(y)$,
ou bien $\Delta(x) \cap \Delta(y) = \emptyset$ et $v(x) \neq v(y)$. On l'a
vérifié sur le modèle ; la page renvoie, comme « justification a priori », au
théorème de la page 2.

\textbf{Bases orthogonales à $\Delta$} (pages 37 et 38). Ce sont les bases
contenues dans $C \smallsetminus \Delta$. Soit $\mathcal{J} = \bigwedge_{i \in
I} \widetilde I_i$ le produit de Baer des $\mathbb{F}_2$-torseurs des deux
permutations circulaires de chaque fibre $T_i$ : un ensemble à deux
éléments\footnote{$\widetilde I_i$ est lu d'après le contexte. Concrètement,
les six transversales se partagent en deux classes de trois selon la parité
de la bijection $T_i \to T_{i'}$ qu'elles définissent, et $\mathcal{J}$ est
l'ensemble de ces deux classes ; c'est l'objet des pages 73 et 74.}. Les bases
orthogonales correspondent aux injections $\varphi : \mathcal{J}
\hookrightarrow V^*$ : $b_\varphi$ est l'ensemble des $x$ dont la classe de
$\Delta(x)$ dans $\mathcal{J}$ s'envoie sur $v(x)$. Il y en a $6$, de sorte que
$(C, \Delta, b)$ équivaut à $\Delta$ seul ($V^* \simeq \mathcal{J} \amalg
\lbrace \mathrm{pt}\rbrace$), et, pour $b$ fixée, il y a $720/72 = 10 =
\tfrac12\binom63$ bitriangles orthogonaux, autant que de partitions de $b$ en
deux moitiés. Il pressent là une « autodualité » de la catégorie des ensembles
à six éléments munis d'une partition de type $(3, 3)$ ; il y revient pages 73
à 75.

\textbf{Les six racines orthogonales à un bitriangle} (page 39). Avec
l'étiquetage des pages 9 et 10 (le bitriangle $u, v, w$, $\alpha_1, \beta_1,
\gamma_1$, $\alpha'_1, \beta'_1, \gamma'_1$), les bases orthogonales sont $b =
(\alpha_1, \alpha'_2, \beta_1, \beta'_2, \gamma_1, \gamma'_2)$ et ses images
par les permutations de $\lbrace 1, 2, 3\rbrace$\footnote{Les indices de la
première ligne se lisent 1 ou 4 ; la suite de la page fait lire 1.} ; la
racine de $b$ est $\alpha_2 - \alpha_1$, et les six racines $\alpha_i -
\alpha_j$ forment un système de type $A_2$.

La page s'achève sur un programme : a) positions d'un bitriangle et d'un
triangle --- trois cas, $t \subset \Delta$ (6), $\operatorname{card} t \cap
\Delta = 1$ ($3 \cdot 9 = 27$), $t \cap \Delta = \emptyset$ (12), total 45 ;
b) positions d'une base et d'un bitriangle, « il semble » trois ; c) « position
relative de deux bitriangles ??? ». Seul a) est fait.

\pagerange{41}{48}
\section*{12. Les vingt-sept droites (pages 41 à 48)}

\subsection*{I. Le plan éclaté (pages 41 à 43)}

Soit $P$ un plan projectif et $S = \lbrace s_1, \ldots, s_6\rbrace$ six points
dont trois ne sont jamais alignés et qui ne sont pas tous sur une conique. Soit
$\widetilde P$ l'éclaté de $P$ en $S$, $D_i$ la courbe exceptionnelle
au-dessus de $s_i$, $D_{ij}$ le transformé strict de la droite $s_i s_j$, $D'_i$
celui de la conique passant par les cinq points autres que $s_i$. Dans
$\operatorname{Pic}(\widetilde P)$, avec $\eta$ l'image inverse de la classe
d'une droite et $\xi_i$ la classe de $D_i$ :
\[
\xi_{ij} = \eta - \xi_i - \xi_j, \qquad \xi'_i = 2\eta - \sum_{\alpha \neq i}
\xi_\alpha = 2\eta - \xi + \xi_i, \qquad \xi = \sum_1^6 \xi_\alpha .
\]
On en tire $\xi' = \sum \xi'_\alpha = 12\eta - 5\xi$, et, définissant $\eta'$
par $2\eta - \xi = -(2\eta' - \xi')$ et posant $\lambda = 3\eta - \xi$,
\[
\xi + \xi' = 4\lambda, \qquad \eta + \eta' = 2\lambda, \qquad \eta' = 5\eta -
2\xi, \qquad \xi_i - \xi'_i = 2\eta' - \xi' .
\]
La forme d'intersection est $\eta^2 = 1$, $\xi_i^2 = -1$, $\eta \cdot \xi_i = 0$,
$\xi_i \cdot \xi_j = 0$ ($i \neq j$), et de même pour $(\eta', (\xi'_i))$. On a
\[
\lambda^2 = 3, \qquad \lambda \cdot \xi_i = \lambda \cdot \xi'_i = \lambda \cdot
\xi_{ij} = 1 ,
\]
\[
\xi_i \cdot \xi'_j = 1 - \delta_{ij}, \qquad \xi_i \cdot \xi_{jk} = \xi'_i \cdot
\xi_{jk} = [i \in \lbrace j, k\rbrace], \qquad \xi_{ij} \cdot \xi_{kl} = 1 -
\operatorname{card}(\lbrace i,j\rbrace \cap \lbrace k, l\rbrace)
\]
pour des paires distinctes. La classe $\lambda$ est celle du système des
cubiques passant par $S$, qui plonge $\widetilde P$ dans $\mathbb{P}^3$ comme
une surface cubique ; $\lambda^2 = 3$ dit que l'image est de degré 3, et
$\lambda \cdot \delta = 1$ que les 27 courbes $D_i, D'_i, D_{ij}$ ($6 + 6 +
15$) deviennent des \emph{droites}. En marge : la classe canonique est
déterminée par $\lambda$, ce qui rend le plongement, et les 27 droites,
intrinsèques à la surface\footnote{La marge écrit « $K = -3\lambda$ ». Pour
l'éclaté de $\mathbb{P}^2$ en six points, $K = -3\eta + \xi = -\lambda$ : la
section hyperplane est anticanonique. La conclusion de la marge (intrinsèque à
$X$) reste juste.}. Chaque droite en rencontre dix autres, et deux droites qui
se rencontrent ont un unique troisième voisin commun : c'est le graphe de la
page 2. Le $(\ast)$ de la page 41, « vérifier », est confirmé par ces
formules, au point près de la transversalité, qu'il note « on espère ! » et
qui résulte de ce que les intersections valent 1.

\subsection*{II. Le même calcul sur un anneau (pages 43 à 48)}

Sur un anneau $k \neq 0$, il part d'un module muni d'éléments $\eta, \xi_i,
\eta', \xi'_i$ liés par (b) $2(\eta + \eta') = \xi + \xi'$ et (c) $\xi'_i -
\xi_i = 2\eta - \xi$, en tire les relations symétriques, $\xi + \xi' = 4\lambda
= 4\lambda'$, et suppose en outre $\lambda = \lambda'$ (conséquence de ce qui
précède si 2 est inversible) ; d'où $2\eta' = 10\eta - 4\xi$, $\eta' = 5\eta -
2\xi$, et symétriquement\footnote{Une formule intermédiaire porte « $3\eta' =
3\eta - \xi + \xi' = 15\eta - 5\xi$ » ; le calcul donne $15\eta - 6\xi$, en
accord avec $\eta' = 5\eta - 2\xi$ qu'il écrit ensuite. Le chiffre repassé est
d'une lecture incertaine.}. La famille $(\eta, (\xi_i))$ est une base si et
seulement si $(\eta', (\xi'_i))$ en est une ; une forme bilinéaire symétrique
vérifiant les relations ci-dessus pour l'une les vérifie pour l'autre, et
elle force la liberté. Les 27 éléments $\xi_i, \xi'_i, \xi_{ij} = \eta - \xi_i -
\xi_j = \eta' - \xi'_i - \xi'_j$ engendrent le module (il suffit des $\xi_i$ et
d'un $\xi_{ij}$).

\textbf{Le graphe standard.} Sur l'ensemble $\Delta$ de ces 27 éléments, on
déclare $\delta \neq \delta'$ liés si $\delta \cdot \delta' = 1$ (les autres
produits sont nuls) : $\xi_i$ est lié aux $\xi'_j$ et $\xi_{ij}$ ($j \neq i$),
$\xi'_i$ aux $\xi_j$ et $\xi_{ij}$, $\xi_{ij}$ à $\xi_i, \xi_j, \xi'_i,
\xi'_j$ et aux six $\xi_{kl}$ disjoints. \textbf{Lemme} (page 46) : deux
sommets liés ont un unique voisin commun --- $\xi_{ij}$ pour $\lbrace \xi_i,
\xi'_j\rbrace$, $\xi'_j$ pour $\lbrace \xi_i, \xi_{ij}\rbrace$, $\xi_j$ pour
$\lbrace\xi'_i, \xi_{ij}\rbrace$, $\xi_{rs}$ pour $\lbrace \xi_{ij},
\xi_{kl}\rbrace$, $\lbrace r, s\rbrace$ complémentaire.

\textbf{Définition.} Un graphe cubique est un graphe isomorphe à ce graphe
standard $C_0$, de sommets $[1,6] \amalg [1,6] \amalg \mathfrak{P}_2([1,6])$.
À $C$ on associe $\widehat E_k(C) = k^{(C)}/N(C)$, quotient par le noyau de la
forme donnée par $\delta \cdot \delta = -1$, $\delta \cdot \delta' = 1$ ou $0$
selon que $\delta, \delta'$ sont liés ou non : un module libre de rang 7 à
forme non dégénérée (c'est $\operatorname{Pic}$, le réseau unimodulaire
$I_{1,6}$).

\textbf{Bases ordonnées} (pages 47 et 48). Une base ordonnée de $C$ est
l'image de $(\xi_1, \ldots, \xi_6)$ par un isomorphisme $C_0 \to C$ ; une
partie basique, l'ensemble sous-jacent. \textbf{Proposition.} $u \mapsto
u((\xi_i))$ est une bijection de $\operatorname{Isom}(C_0, C)$ sur les bases
ordonnées, qui forment donc un torseur, à droite sous $W_0 =
\operatorname{Aut}(C_0)$ et à gauche sous $W$ ; et les automorphismes qui
stabilisent une partie basique sont induits par $\mathfrak{S}_6$. Preuve :
$\xi'_i$ est l'unique sommet lié à tous les $\xi_j$, $j \neq i$, et $\xi_{ij}$
l'unique sommet lié à $\xi_i$ et $\xi'_j$ ; un automorphisme qui fixe les
$\xi_i$ est donc l'identité. Une base $(\alpha_i)_{i \in I}$ équivaut à un
isomorphisme $C \simeq I \amalg I \amalg \mathfrak{P}_2(I)$.

\textbf{Proposition.} Il existe un unique $\lambda_C \in \widehat E(C)$ tel que
$\lambda \cdot s = 1$ pour tout sommet $s$ ; pour toute base $\alpha$,
\[
\lambda = 3\eta(\alpha) - \xi(\alpha), \quad \eta(\alpha) + \eta(\alpha') =
2\lambda, \quad \xi(\alpha) + \xi(\alpha') = 4\lambda, \quad \sum_{ij}
\xi_{ij}(\alpha) = 5\lambda, \quad \sum_{s \in C} s = 9\lambda .
\]
L'unicité vient de ce que les sommets engendrent et que la forme est non
dégénérée ; l'existence se voit sur $C_0$. $\lambda$ est donc invariant par
automorphisme, et le produit scalaire avec $\lambda$ donne une suite exacte
invariante $0 \to E \to \widehat E \to \mathbb{Z} \to 0$, $E = \lambda^\perp$,
de rang 6.

\pagerange{49}{57}
\section*{13. Relations avec $E_6$ (pages 49 à 57)}

\subsection*{Racines et double-six (pages 49 à 51)}

\textbf{Définition.} Une \emph{racine} de $C$ est une différence $\delta' -
\delta$ de sommets distincts non liés ; leur ensemble $R(C) \subset E(C)$ est
stable par $r \mapsto -r$, et $r^2 = -2$.

Les $27 \cdot 16 = 432$ couples ordonnés de sommets distincts non liés donnent
72 racines, six couples par racine : pour chaque racine $r$, il y a un
double-six $\alpha = (\alpha_1, \ldots, \alpha_6)$, $\alpha' = (\alpha'_1,
\ldots, \alpha'_6)$, avec $\alpha_i$ lié à $\alpha'_j$ si et seulement si $i
\neq j$, les $\alpha_i$ deux à deux non liés, de même les $\alpha'_i$, et $r =
\alpha'_i - \alpha_i$ pour tout $i$. Dans $C_0$ :
\[
\begin{array}{lll}
1 & r_0 = 2\eta - \xi = \lambda - \eta = \xi'_i - \xi_i &
(\xi_1, \ldots, \xi_6), \ (\xi'_1, \ldots, \xi'_6) \\[3pt]
15 & r_{ij} = \xi_i - \xi_j \ (i < j) & (\xi_j, \xi'_j, (\xi_{ik})_{k \neq
i,j}), \ (\xi_i, \xi'_i, (\xi_{jk})_{k \neq i,j}) \\[3pt]
20 & r_{ijk} = \eta - \xi_i - \xi_j - \xi_k & (\xi_i, \xi_j, \xi_k,
\xi_{mn}, \xi_{ln}, \xi_{lm}), \ (\xi_{jk}, \xi_{ik}, \xi_{ij}, \xi'_l,
\xi'_m, \xi'_n)
\end{array}
\]
avec $\lbrace l, m, n\rbrace$ le complémentaire de $\lbrace i, j, k\rbrace$, et
les opposées, obtenues en échangeant $\alpha$ et $\alpha'$ : $2 \cdot (1 + 15
+ 20) = 72$\footnote{Pour $r_{ij}$, la transcription note que les indices
« se lisent ainsi, bien que la définition des sixtes demande $(\xi_j, \xi'_i,
\ldots)$ » ; c'est la page qui a raison : $(\xi_j, \xi'_j, (\xi_{ik}))$ et
$(\xi_i, \xi'_i, (\xi_{jk}))$ forment bien un double-six de racine $\xi_i -
\xi_j$ (vérifié), et la page 59 donne la même liste. En revanche
$\alpha^{ij} = (\xi_i, \xi'_j, (\xi_{ik}))$, écrit page 51, et $(\xi_j,
\xi'_i, (\xi_{ik}))$, page 50, ne sont pas des bases ($\xi_i$ est lié à
$\xi_{ik}$). Pour $r_{ijk}$ on a rangé la première sixte pour que $\alpha'_i -
\alpha_i = r$ terme à terme ; la page écrit $(\xi_{lm}, \xi_{mn}, \xi_{nl})$,
même ensemble. Ce sont les double-six de Schläfli : 36 double-six, 72
sixains.}.

Pour une racine $r$ de double-six $(\alpha, \alpha')$, on a $r \cdot \alpha_i =
1$, $r \cdot \alpha'_i = -1$, et $r \cdot \delta = 0$ pour les quinze autres
sommets ; la racine détermine donc ses deux sixains comme $\lbrace s : r \cdot s
= 1\rbrace$ et $\lbrace s : r \cdot s = -1\rbrace$, et les racines sont en
bijection avec les bases. Pour $r$ donnée, $\eta(\alpha) = \lambda - r$ et
$\xi(\alpha) = \sum \alpha_i = 2\lambda - 3r$ (passage biffé de la page 50).

\subsection*{Réflexions et groupe de Weyl (pages 51 à 53)}

Après « Vérifications fastidieuses\ldots », un « meilleur procédé » : comme $r
\cdot r = -2$,
\[
s_r(x) = x + (x \cdot r)\, r
\]
est la réflexion orthogonale par rapport à $r^\perp$ ; elle préserve la forme
et $\lambda$, et $s_r \alpha_i = \alpha'_i$, $s_r \alpha'_i = \alpha_i$, $s_r
\delta = \delta$ pour les autres sommets : $s_r$ est un automorphisme de $C$.
Soit $W$ le groupe engendré par les $s_r$. $s_{r_{ij}}$ est la transposition
$(i\, j)$, donc $W \supset \mathfrak{S}_6$ ; et
\[
s_{r_{ijk}}(r_0) = r_{lmn}, \qquad s_{r_{ijk}}(r_{il}) = r_{jkl},
\]
de sorte que $W$ est transitif sur les racines. \textbf{Corollaires} : $W$ est
transitif sur les racines, les sixains, les bases ordonnées, les couples de
sommets non liés, et $W$ est le groupe de \emph{tous} les automorphismes de
$C$ ; simplement transitif sur les $72 \cdot 6!$ bases ordonnées, il est
d'ordre
\[
\operatorname{card} W = 72 \cdot 6! = 2^7 \cdot 3^4 \cdot 5 .
\]

\subsection*{Le théorème (pages 53 à 55)}

\textbf{Théorème.} $(E(C), R(C))$ est un système de racines de type $E_6$, et
si $(\xi_1, \ldots, \xi_6)$ est une base ordonnée,
\[
r_1 = \xi_1 - \xi_2, \ r_2 = \xi_2 - \xi_3, \ r_3 = \xi_3 - \xi_4, \ r_4 =
\xi_4 - \xi_5, \ r_5 = \xi_5 - \xi_6, \ r_6 = \eta - \xi_1 - \xi_2 - \xi_3
\]
est un système de racines simples, de diagramme

\begin{tikzcd}[column sep=small, row sep=small]
r_1 \arrow[r, no head] & r_2 \arrow[r, no head] & r_3 \arrow[r, no head] \arrow[d, no head] & r_4 \arrow[r, no head] & r_5 \\
& & r_6 & &
\end{tikzcd}

$W$ s'identifie au groupe de Weyl, et la forme est la forme invariante
normalisée par $r \cdot r = -2$.

La démonstration : $E$ est engendré par $r_1, \ldots, r_6$ (les $\xi_i - \xi_1$
et $\eta - 3\xi_1$ engendrent $E$, et $\eta - 3\xi_1 = r_6 - (\xi_1 - \xi_2) -
(\xi_1 - \xi_3)$) ; les racines sont stables par les $s_r$ et deux racines non
opposées ne sont pas proportionnelles ; puis toute racine positive est
combinaison à coefficients entiers positifs des $r_i$ :
\[
r_{ij} = r_i + \cdots + r_{j-1}, \qquad r_{ijk} = r_6 + (r_1 + \cdots +
r_{i-1}) + (r_2 + \cdots + r_{j-1}) + (r_3 + \cdots + r_{k-1}),
\]
de hauteur $i + j + k - 5 \leq 10$, et
\[
r_0 = r_{123} + r_{456} = r_1 + 2r_2 + 3r_3 + 2r_4 + r_5 + 2r_6 ,
\]
de hauteur 11 : c'est la plus grande racine. Enfin $r_i \cdot r_j = 0$ sauf
pour $j = i + 1 \leq 5$ et pour $(i, j) = (3, 6)$, où il vaut 1 : le diagramme
est $E_6$\footnote{Les coefficients de $r_1, r_2, r_3$ dans l'expression de
$r_0$ sont surchargés et le « 11 » est écrit sur un autre nombre ; on a
vérifié $r_0 = (1, 2, 3, 2, 1, 2)$ dans la base des $r_i$, la plus grande
racine de $E_6$.}.

\textbf{Corollaire.} $W = \operatorname{Aut}(C)$ est d'indice 2 dans le groupe
$\widetilde W$ des automorphismes du système de racines, le quotient étant le
groupe des automorphismes du diagramme de Dynkin.

\subsection*{Retrouver le graphe à partir des racines (pages 56 et 57)}

La donnée de $C$ est donc un peu plus fine que celle du système de racines :
elle revient au choix de l'une des deux extrémités longues du diagramme. Il
voudrait voir, dans la géométrie interne de $(E, R)$, ce qui récupère $C$. Il
pose
\[
i(\delta) = 3\delta - \lambda \in E \qquad (\delta \cdot \lambda = 1, \ \lambda^2 = 3),
\]
d'où $i(\delta) \cdot i(\delta') = 9\, \delta\cdot\delta' - 3$ :
\[
i(\delta)^2 = -12, \qquad i(\delta) \cdot i(\delta') = \begin{cases} 6 &
\delta, \delta' \text{ liés} \\ -3 & \delta \neq \delta' \text{ non
liés.}\end{cases}
\]
$i$ est une bijection de $C$ sur une partie $\Delta'$ de l'ensemble des
vecteurs de carré $-12$ de $E$, invariante par $W$, et les racines s'en
déduisent par\footnote{La page écrit « $9$ si $\delta, \delta'$ liés » et « $x, x'$ liés
$\iff x \cdot x' = 9$ » ; le calcul donne $9 - 3 = 6$. Elle annonce
« (72) » vecteurs de carré $-12$, lecture incertaine : il y en a bien plus
(2214 dans $E$) ; la partie $\Delta' \cup (-\Delta')$ en compte 54. Enfin elle
écrit $3R = \lbrace x' - x\rbrace$ sans restriction ; pour $x, x'$ liés,
$x' - x$ est trois fois un vecteur de carré $-4$, qui n'est pas une
racine.}
\[
3R = \lbrace x' - x \mid x, x' \in \Delta', \ x \cdot x' = -3\rbrace
\]


$-\Delta'$ est aussi invariant par $W$, et $\Delta' \cap (-\Delta') =
\emptyset$, car $3(\delta + \delta') = 2\lambda$ est impossible ($2\lambda =
6\eta - 2\xi$ n'est pas divisible par 3). Donc $-\mathrm{id}_E \notin W$ et\footnote{La page écrit « engendré par
$\operatorname{id}_E$ », pour $-\operatorname{id}_E$.}
\[
\widetilde W \simeq W \times \lbrace \pm \mathrm{id}_E\rbrace .
\]
La réunion $\widetilde\Delta = \Delta' \cup
(-\Delta')$ est canoniquement attachée au système de racines, se décompose en
deux orbites sous $W$, et le choix d'une orbite équivaut au choix de $C$. « On
aimerait l'expliciter, ainsi que $\widetilde\Delta$, via la géométrie interne
du syst.\ de racines. » Ce sera l'objet des pages 95 à 98 ; la réponse est
que $\tfrac13 \widetilde\Delta$ est l'ensemble des \emph{poids minuscules} non
nuls de $E_6$ (section 21).

\pagerange{58}{60}
\section*{14. Une définition intrinsèque, et les 72 bases (pages 58 à 60)}

\textbf{Définition} (page 58). Un graphe $C$ est \emph{cubique} s'il existe
deux parties disjointes $I$, $I'$ de six sommets, chacune formée de sommets
deux à deux non liés\footnote{Cette condition n'est pas écrite et elle est
nécessaire : sans elle rien ne fixe les liaisons à l'intérieur de $I$ et de
$I'$. On l'ajoute.}, telles que

\begin{enumerate}
\item[b)] tout $i \in I$ est non lié à un unique $i' \in I'$, et
  inversement\footnote{La page écrit deux fois « non lié à $I$ » pour « non
  lié à $i$ », « à $i'$ ».}, d'où une bijection $i : I \to I'$ ;
\item[c)] pour toute paire $J \subset I$, il existe un unique sommet $\xi_J$
  lié aux quatre éléments de $J \cup i(J)$ et à aucun autre élément de $I \cup
  I'$ ;
\item[d)] $J \mapsto \xi_J$ est une bijection de $\mathfrak{P}_2(I)$ sur $C
  \smallsetminus (I \cup I')$ ;
\item[e)] $\xi_J$ et $\xi_{J'}$ ($J \neq J'$) sont liés si et seulement si $J
  \cap J' = \emptyset$.
\end{enumerate}

Avec l'hypothèse ajoutée, ces conditions décrivent exactement le graphe
$C(I)$ construit sur un ensemble $I$ à six éléments, et $C$ est cubique si et
seulement si $C \simeq C(I)$. $I$ détermine $I'$ ; ce sont les \emph{parties
basiques} ; une \emph{bibase} est une paire de bases associées.

\textbf{Proposition.} a) La catégorie des graphes cubiques munis d'une base
(resp.\ d'une base ordonnée) équivaut à celle des ensembles à six éléments
(resp.\ à la catégorie ponctuelle) ; b) munis d'une bibase, à celle des couples
d'un ensemble à six éléments et d'un ensemble à deux éléments ; c) les bases
ordonnées forment un torseur sous $W_0$ à droite et $W$ à gauche.

\textbf{Théorème} (pages 59 et 60). Les bases de $C(I)$ sont : (i) $I$ et $I'$
(2) ; (ii) pour $(i, j) \in I^2$, $i \neq j$, $B^{(i,j)} = \lbrace \xi_j,
\xi'_j, (\xi_{ik})_{k \neq i, j}\rbrace$ (30) ; (iii) pour $K \in
\mathfrak{P}_3(I)$, $B^K = \lbrace \xi_i \ (i \in K), \ \xi_J \ (J \in
\mathfrak{P}_2(I \smallsetminus K))\rbrace$ (20) ; (iv) $B'^K = \lbrace \xi'_i \ (i \notin K), \ \xi_J \ (J \in \mathfrak{P}_2(K))\rbrace$ (20). La base associée
à $B^{(i,j)}$ est $B^{(j,i)}$, celle de $B^K$ est $B'^K$. \textbf{Corollaires} :
il y a 72 bases et $\operatorname{card} W = 72 \cdot 6!$ ; $W$ est transitif
sur les sommets et sur les couples de sommets distincts non liés, et pour un
tel couple $(\delta, \delta')$ il existe une unique base $I$ contenant $\delta$
dont la base associée contient $\delta'$ ; la catégorie des graphes cubiques
munis d'une paire de sommets non liés équivaut à celle des couples d'un
ensemble à cinq éléments et d'un ensemble à deux éléments.

La fin de la page 60, barrée de deux croix, ébauche une notion de
« bi-base » et étudie les parties de trois sommets liés ; elle ne se lit
qu'en partie et on ne la restitue pas.

\pagerange{61}{63}
\section*{15. Un tableau des types, et les produits de racines (pages 61 à 63)}

La page 61 est un tableau : pour chaque type de configuration, la catégorie
équivalente, le groupe d'automorphismes, le nombre de configurations et les
longueurs des orbites du stabilisateur sur les 27 sommets. On le redonne
corrigé, chaque ligne vérifiée sur le modèle :
\[
\begin{array}{llrrl}
\text{type} & \text{catégorie} & \text{groupe} & \text{nombre} &
\text{orbites} \\ \hline
1\text{-base} & Q \text{ (5) et torseur} & 1920 & 27 & 1, 10, 16 \\
2\text{-base} & \mathrm{Ens}_5 \times \mathrm{Ens}_2 & 240 & 216 & 2, 5, 10, 10
\\
2\text{-base ordonnée} & \mathrm{Ens}_5 & 120 & 432 & 1, 1, 5, 5, 5, 10 \\
3\text{-base} & \mathrm{Ens}_3 \times \mathrm{Ens}_3 \times \mathrm{Ens}_2 &
72 & 720 & 3, 3, 6, 6, 9 \\
3\text{-base ordonnée} & \mathrm{Ens}_3 \times \mathrm{Ens}_2 & 12 & 4320 &
1^3, 2^3, 3^4, 6 \\
4\text{-base} & \mathrm{Ens}_4 \times \mathrm{Ens}_2 & 48 & 1080 & 1, 2, 2,
4, 4, 6, 8 \\
4\text{-base ordonnée} & \mathrm{Ens}_2 & 2 & 25\,920 & 1^{15}, 2^6 \\
5\text{-base} & \mathrm{Ens}_5 & 120 & 432 & 1, 1, 5, 5, 5, 10 \\
5\text{-base ordonnée} & \text{ponctuelle} & 1 & 51\,840 & 1^{27} \\
6\text{-base} & \mathrm{Ens}_6 & 720 & 72 & 6, 6, 15 \\
\text{demi-étoile} & \mathrm{Ens}_5 \times \mathrm{Ens}_2 & 240 & 216 & 2, 5,
10, 10 \\
\text{triades} & & 1296 & 40 & 27 \\
\text{triades ordonnées} & & 216 & 240 & 9, 9, 9 \\
\text{grilles} & & 432 & 120 & 9, 18
\end{array}
\]
Ici une $n$-base est une partie de $n$ sommets deux à deux non liés contenue
dans une base, une demi-étoile la partie $E(x) \cap E(y)$ pour $x, y$ non liés
(page 66), une grille un bitriangle, une triade une partition des 27 sommets en
trois bitriangles\footnote{Différences avec la page : pour la 2-base, la page
donne les orbites $1, 5, 10, 10$ (somme 26 ; les deux sommets de la paire
forment une orbite de longueur 2) ; pour la 2-base ordonnée $1, 5, 5, 5, 10$
(il manque un 1) ; pour la 3-base ordonnée la suite, écrite serré, se lit
$1,1,1,2,2,2,2,3,3,3,3,6$ (somme 29) ; pour la 4-base $1, 2, 3, 4, 4, 6, 8$ ;
pour la demi-étoile, le nombre se lit « 15 » et les orbites renvoient à celles
de la 2-base ordonnée. Les 40 triades sont les 40 façons de partager les 27
droites en trois paires trièdres (nom de nous).}.

\textbf{Page 63.} Les racines dans la base $\eta, \xi_1, \ldots, \xi_6$, avec
leurs double-six ($r_{ijk} = \xi_{jk} - \xi_i = \xi'_l - \xi_{mn} = \cdots$),
et la table des produits :
\[
r_0 \cdot r_{ij} = 0, \quad r_0 \cdot r_{ijk} = -1, \quad r_{ijk} \cdot
r_{pqs} = 1 - \operatorname{card}(\lbrace i,j,k\rbrace \cap \lbrace
p,q,s\rbrace),
\]
$r_{ij} \cdot r_{i'j'} = 0$ si $\lbrace i,j\rbrace \cap \lbrace i',j'\rbrace =
\emptyset$, $+1$ si $i' = j$ ou $i = j'$, $-1$ si $i = i'$ ou $j =
j'$\footnote{La page porte « $-1$ si $i \neq i'$ ou $j = j'$ » ; le signe est
peut-être un $=$ mal formé.} ; $r_{ij} \cdot r_{pqs}$ vaut $0$ si $\lbrace i,
j\rbrace$ est contenu dans $\lbrace p, q, s\rbrace$ ou lui est disjoint, $+1$ si
seul $i$ y est, $-1$ si seul $j$ y est. Il en tire $\cos\widehat{(r, r')} = -
r \cdot r'/2$, et dessine l'étoile $A_2$ de $r_0$, $r_{123} = r_6$,
$r_{456}$, avec\footnote{La page écrit $r_0 = 3(r_1 + r_2 + r_3) + 2r_4 + r_5 + 2r_6$ et
$r_{456} = r_6 + 3(r_1 + r_2 + r_3) + 2r_4 + r_5$, l'écriture laissant voir
des $r$ récrits sur des $\xi$ ; ces coefficients sont ceux de $3\xi_1 - \xi_4
- \xi_5 - \xi_6$, non des racines. La page 55 a la bonne expression de $r_0$.
La page 62, tête-bêche, est le brouillon de ce calcul.}
\[
r_{456} = r_1 + 2r_2 + 3r_3 + 2r_4 + r_5 + r_6, \qquad r_0 = r_6 + r_{456} .
\]

\pagerange{66}{71}
\section*{16. Les parties libres (pages 66 à 71)}

Une partie \emph{$n$-libre} est une partie de $n$ sommets deux à deux non
liés.

\textbf{Proposition} (page 66). (1) Il existe des parties $n$-libres si et
seulement si $1 \leq n \leq 6$ ; les parties libres maximales sont les bases
(cardinal 6) et les \emph{deux-étoiles} $E(x) \cap E(y)$, $x, y$ distincts non
liés (cardinal 5). (2) Soit $L_i$ le graphe discret à $i$ sommets ;
$\operatorname{Pl}(L_i, C)$ est homogène sous $W$ pour $i \neq 5$, et a deux
orbites pour $i = 5$ : les 5-bases et les deux-étoiles. Les restrictions
\[
\operatorname{Pl}(L_6) \xrightarrow{\ 1\ } \operatorname{Pl}'(L_5)
\xrightarrow{\ 2\ } \operatorname{Pl}(L_4) \xrightarrow{\ 6\ }
\operatorname{Pl}(L_3) \xrightarrow{\ 10\ } \operatorname{Pl}(L_2)
\xrightarrow{\ 16\ } \operatorname{Pl}(L_1) = C \xrightarrow{\ 27\ }
\mathrm{pt},
\]
et $\operatorname{Pl}''(L_5) \to \operatorname{Pl}(L_4)$ de degré 1, sont
surjectives de degrés indiqués. D'où les nombres de parties libres\footnote{La page écrit « $\simeq B''_4(C)$ » sous $B_4$ : il s'agit des
deux-étoiles, $B''_5 = 25\,920/5! = 216$, en bijection avec les paires de
sommets non liés. On a vérifié tous ces nombres, et que les parties 5-libres
maximales sont exactement les 216 deux-étoiles.} :
\[
\begin{gathered}
B_1 = 27, \ B_2 = 216, \ B_3 = 720, \ B_4 = 1080, \ B'_5 = 432, \ B''_5 = 216,
\ B_6 = 72, \\ \operatorname{card} \operatorname{Bib}(C) = 36 .
\end{gathered}
\]

\textbf{Étude cas par cas.} 1) Le stabilisateur d'un sommet est transitif sur
ses 16 non-voisins. 2) Pour un couple $(x, y)$ de sommets non liés, il existe
une unique base $b$ avec $x \in b$, $y \in b'$ ; avec $(x, y) = (\xi_6, \xi'_6)$,
les sommets non liés à $x$ ni à $y$ sont les dix $\xi_{ij}$, $i, j \leq 5$, sur
lesquels $\mathfrak{S}_5$ est transitif. Une paire non liée équivaut à une
bibase avec un point de $I = B/\sigma_B$ (page 68).

3) Une partie 3-libre a le type $J = \lbrace \xi_{12}, \xi_{23},
\xi_{31}\rbrace$ ; ses non-voisins sont $\xi_4, \xi_5, \xi_6, \xi'_4, \xi'_5,
\xi'_6$. Les parties 3-libres correspondent aux couples $(B, K)$ d'une bibase
et d'une partie à trois éléments de $I$, modulo $K \sim I \smallsetminus K$ (72
$\cdot 20/2 = 720$) ; la catégorie des $(C, J)$ équivaut à celle des triples
$(\tau, K, K')$ d'ensembles à 2, 3, 3 éléments, avec $C = (\tau \times (K
\amalg K')) \amalg \mathfrak{P}_2(K \amalg K')$, d'où un groupe
$\mathfrak{S}_2 \times \mathfrak{S}_3 \times \mathfrak{S}_3$ d'ordre 72.
L'involution $(\tau, K, K') \mapsto (\tau, K', K)$ associe à $J$ l'ensemble $J'$
des sommets liés à tout $J$ ($\lbrace \xi_{45}, \xi_{56}, \xi_{64}\rbrace$).
\textbf{Corollaire} : $J$ est contenu dans exactement deux bases, dont
l'intersection est $J$ ; les parties 3-libres correspondent donc aux paires de
bases d'intersection 3, et aux paires de racines à $60^\circ$ (page 89).

\textbf{Corollaire} (page 70). Pour $J$ 3-libre, l'ensemble $\mathcal{L}(J)$
des sommets non liés à $J$ est un \emph{hexagone} (un 6-cycle induit), et $J
\mapsto \mathcal{L}(J)$ est une bijection des 720 parties 3-libres sur les 720
hexagones de $C$, d'inverse $\mathcal{H} \mapsto \mathcal{L}(\mathcal{H})$. Un
hexagone équivaut à un couple $(\tau, K')$ de cardinaux 2 et 3 (le graphe
$\tau \times K'$, deux sommets liés si leurs deux composantes diffèrent).

4) Pour $J = \lbrace \xi_1, \ldots, \xi_4\rbrace$, $\mathcal{L}(J) = \lbrace
\xi_5, \xi_6, \xi_{56}\rbrace$, en deux orbites ($\lbrace \xi_5, \xi_6\rbrace =
b \cap \mathcal{L}(J)$ et le point $\xi_{56}$) : d'où les degrés 2 et 1 au-dessus
de $\operatorname{Pl}(L_4)$ ; l'ensemble $E(J)$ des sommets liés à tout $J$ est
$\lbrace \xi'_5, \xi'_6\rbrace$, en bijection avec $b \smallsetminus J$ par
$\sigma_b$. 5) Une partie 5-libre est contenue dans une base, alors unique, ou
est une deux-étoile ; $\lbrace x, y\rbrace \mapsto E(x) \cap E(y)$ est une
bijection des paires non liées sur les deux-étoiles, d'inverse $J \mapsto
E(J)$ ; pour $z = \lbrace \xi_6, \xi'_6\rbrace$, $J = \lbrace \xi_{i6}\rbrace_{i
\leq 5}$.

\textbf{Bases et bibases} (« il fallait commencer par là ! »). $W$ est
simplement transitif sur les bases ordonnées ; 72 bases de stabilisateur
$\mathfrak{S}_6$, 36 bibases de stabilisateur $\mathfrak{S}_2 \times
\mathfrak{S}_6$.

\pagerange{73}{75}
\section*{17. Le bitriangle est auto-complémentaire (pages 73 à 75)}

1) Un \emph{bitriangle} est un graphe isomorphe au produit $S \times T$ de
deux triangles : $(\alpha, \beta)$ et $(\alpha', \beta')$ distincts sont liés
si et seulement si $\alpha = \alpha'$ ou $\beta = \beta'$. Neuf sommets, 18
arêtes ; ses six triangles sont les fibres des deux projections, et se
groupent en deux familles (deux triangles d'une même famille sont disjoints),
d'où un revêtement de degré 3 $T \to \tau$ sur un ensemble à deux éléments.
La catégorie des bitriangles équivaut à celle des ensembles à six éléments
munis d'une partition de type $(3,3)$.

2) Le \emph{bloc de triades} est le graphe $D = S \times T$ où $(\lambda, \mu)$
et $(\lambda', \mu')$ sont liés si et seulement si $\lambda \neq \lambda'$ et $\mu
\neq \mu'$ : le complémentaire du bitriangle, avec lui aussi 18 arêtes. Ses
six triades (parties de trois sommets deux à deux non liés) se groupent en
deux familles de trois.

\textbf{Proposition.} Le bloc de triades est isomorphe à un bitriangle :
$\Delta \simeq \overline\Delta$.

\textbf{Démonstration.} Les triades du bitriangle $S \times T$ sont les
graphes des six bijections $S \to T$, et forment un torseur sous
$\mathfrak{S}_S$ ; on les partage en deux classes selon la parité de
$\varphi^{-1}\psi$. Deux bijections distinctes d'une même classe diffèrent par
une permutation circulaire, sans point fixe : leurs graphes sont disjoints, et
chaque classe est une partition des neuf sommets en trois. D'où une
application $\Delta \to \prod_{i \in \sigma} \mathfrak{S}_i$ (la classe dans
chaque famille), qui est bijective : une fois $S$ identifié à $T$, un couple
$(\alpha, \beta)$ est déterminé par la permutation circulaire $\rho_1$ telle
que $\beta = \rho_1 \alpha$ et la transposition $\theta$ telle que $\beta =
\rho_1 \theta\alpha$, dont $\alpha$ est l'unique point fixe. Deux sommets liés
dans $\Delta$ ont des images distinctes dans les deux familles, donc sont liés
pour la structure de bloc de triades transportée par cette bijection ; les
deux structures ont 18 arêtes, elles coïncident.

\textbf{Conclusion.} $\Delta \mapsto \overline\Delta$ est une autoéquivalence
involutive de la catégorie des bitriangles, donc de celle des ensembles à six
éléments munis d'une partition de type $(3, 3)$. La page 75 commence à définir
une « dualité » entre deux tels objets $T$ et $S$ --- une bijection $\prod T
\simeq \prod S$ telle que\ldots --- et s'arrête\footnote{C'est
l'« autodualité » annoncée page 38. Dans le langage d'aujourd'hui : le
bitriangle est le graphe de Paley d'ordre 9, auto-complémentaire, et la
démonstration revient à voir $S \times T$ comme le plan affine
$\mathrm{AG}(2, 3)$, dont les quatre directions sont les lignes, les colonnes
et les deux classes de bijections ; le complémentaire échange la paire
« lignes, colonnes » et la paire « diagonales ». Noms de nous.}.

\pagerange{77}{80}
\section*{18. Polyèdres et polygones dans le graphe (pages 77 à 80)}

Pages de dénombrements et de dessins. Polygones (sous-graphes induits) :
\[
\text{carrés } 27 \cdot 10 \cdot 4 = 1080, \qquad \text{pentagones } 36 \cdot 6
\cdot 12 = 2592, \qquad \text{hexagones } 720 ,
\]
nombres qu'on a vérifiés ; polyèdres : prismes $45 \cdot 16 = 720$, prismes
doubles $45 \cdot 8$, antiprismes tordus, « pseudocubes », cubes $27 \cdot 10
\cdot 2$, dodécaèdres gauches $2^4 \cdot 3^3 = 432$, avec un dodécaèdre gauche
dessiné et étiqueté par des $\xi_{ij}$\footnote{Pour les antiprismes, la page
écrit $45 \cdot 8 \cdot 16 \cdot 3 = 2^4 3^4 5$ ; le produit vaut $17\,280 =
2^7 3^3 5$, et le « 8 » est récrit sur un autre chiffre. On n'a pas vérifié
ces comptes de polyèdres, dont les définitions ne sont pas données.}. Au bas
de la page 77, tête-bêche, un résumé de la construction $\mathbb{Z}^{(C)} \to
\widehat E(C) \supset E(C) = \lambda^\perp$ de la page 46. La page 78 reprend
en brouillon les sections de la page 17 ($M \subset \mathbb{F}_2^Q$, $Q$ à
cinq éléments). La page 79 donne, entre autres, un groupe d'ordre 144, $(
\mathfrak{S}_3 \times \mathfrak{S}_3) \cdot \mathfrak{S}_2 \times
\mathfrak{S}_2$, et décrit des « pentaprismes » comme des couples de 3-bases
munis d'une bijection, c'est-à-dire une bibase et une partition de type $(3,
3)$ avec un isomorphisme des deux moitiés ; en bas, une question sur trois
automorphismes $U, V, W$ fixant $P_u, P_v, P_w$ : « $UVW \in \lbrace 1,
\sigma\rbrace \Rightarrow$? $U, V, W = 1$ », et de nouveau l'écriture $I \times
\lbrace 0, 1, 2\rbrace \amalg I \times J \times \operatorname{circ}(I)$ de la
page 13. La page 80 énonce un lemme sur un triangle et deux triangles voisins,
dont l'implication est barrée ; il reste une question.

\pagerange{81}{87}
\section*{19. Différences, sommes, indices (pages 81 à 87)}

La page 81 est un dessin : un fragment à faces pentagonales, étiqueté par des
$\xi_{ij}$. La page 82 classe les différences et les sommes de sommets :

\begin{itemize}
\item[(0)] les différences de sommets non liés se groupent par six parmi les
  $27 \cdot 16$ couples : les 72 racines ;
\item[(6)] les différences de sommets liés sont 270 éléments distincts, de
  carré $-4$, d'où $72 + 270 = 342$ éléments\footnote{On lit aussi « 392 » ;
  $72 + 270 = 342$.} ;
\item[(1)] les sommes de sommets \emph{liés} se groupent par cinq parmi les
  135 paires : $\delta + \delta' = \lambda - \delta''$, $\delta''$ le troisième
  sommet du triangle, d'où 27 éléments, de carré $0$\footnote{La page écrit
  « sommets non liés » et « de carré 2 ». Le nombre 135 est celui des paires
  liées, et $(\lambda - \delta'')^2 = 3 - 2 - 1 = 0$.} ;
\item[(b)] les sommes de sommets non liés donnent $216$ éléments distincts,
  de carré $-2$, soit $27 + 216 = 243$ en tout.
\end{itemize}

La page 83 est l'annonce dactylographiée décrite plus haut. Les pages 84 et 86
calculent des indices dans $W$ de stabilisateurs de configurations épinglées
ou pointées et orientées --- prismes ($[W] : 6$, puis 72), hexagones
($8640$, puis 720), pentagones ($25\,920$, puis 2592), carrés ($8640$, puis
1080), triangles pointés orientés ($270$), triangles ($45 = [W] : 1152$),
« carrousels » ($[W] : 2$, puis 12) --- et une chaîne d'équivalences
« $A_4$ ép.\ $\sim A_5$ ép.\ $\sim$ hex.\ ép.\ $\sim$ prisme épinglé $\sim$
carrous.\ ép. », que la page n'explique pas\footnote{Les noms $A_3$, $A_4$,
$A_5$ y désignent des configurations dessinées, non définies ; on ne les
interprète pas. Pour les triangles pointés orientés, la page donne l'indice
avec des exposants repassés ; c'est $51\,840/270 = 192$.}.

\textbf{Un codage par parité} (page 86). Les sommets voisins d'une étoile
$A$ sont codés par des vecteurs de $\mathbb{F}_2^A$, et la liaison par le poids
de leur différence : « $x, y$ liés $\iff \operatorname{card} u = 3$ ou $4$ »,
« non liés $\iff \operatorname{card} u = 0, 1$ ou $2$ », $u = x - y$. Sur des
vecteurs de même parité, c'est la règle d) de la page 17 (poids 4 pour les
liés, 2 pour les non liés). La page écrit aussi, pour la même parité, « liés
$\iff x - y = \mathbf{1}^A$ », ce qui est impossible si $A$ a cinq éléments ;
le reste du codage ne se laisse pas reconstituer avec certitude. On note
aussi la formule $\delta \cdot \delta' \equiv \operatorname{card}(\delta \cap
\delta') + \operatorname{card}(\varepsilon(\delta) \cap \varepsilon(\delta'))$,
sans définition de $\varepsilon$. La page 87 compte les « bicarrousels » :
$1080 = 72 \cdot 15$, d'indice 48.

\pagerange{89}{92}
\section*{20. Positions de deux racines, et les 27 en coordonnées (pages 89 à
92)}

\textbf{Positions relatives de deux racines} (page 89). Pour des racines $r,
r'$ de bases $b, b'$ :
\[
\begin{array}{llll}
r' = r & r \cdot r' = -2 & b = b' & 1 \\
r' = -r & r \cdot r' = 2 & b' \text{ associée à } b & 1 \\
r \perp r' & r \cdot r' = 0 & \operatorname{card}(b \cap b') = 1 & 30 \\
\widehat{(r, r')} = 60^\circ & r \cdot r' = -1 & \operatorname{card}(b \cap b')
= 3 & 20 \\
\widehat{(r, r')} = 120^\circ & r \cdot r' = 1 & \operatorname{card}(b \cap b')
= 0, \ b, b' \text{ non associées} & 20
\end{array}
\]
(vérifié). D'où $72 \cdot 15 = 1080$ paires orthogonales, $720$ paires à
$60^\circ$, $720$ à $120^\circ$, et $720/6 = 120$ « plans spéciaux », les plans
engendrés par deux racines non orthogonales : les 120 sous-systèmes de type
$A_2$.

\textbf{Les 27 dans la base des racines simples} (pages 90 à 92). Il exprime
$\widetilde\xi_i = 3\xi_i - \lambda$ (le $i(\delta)$ de la page 56) dans la
base $r_1, \ldots, r_6$ ; la table de la page 92 est juste pour les $\xi_i$ :
\[
\begin{array}{c|cccccc}
& r_1 & r_2 & r_3 & r_4 & r_5 & r_6 \\ \hline
-\widetilde\xi_1 & -1 & 1 & 3 & 2 & 1 & 3 \\
-\widetilde\xi_2 & 2 & 1 & 3 & 2 & 1 & 3 \\
-\widetilde\xi_3 & 2 & 4 & 3 & 2 & 1 & 3 \\
-\widetilde\xi_4 & 2 & 4 & 6 & 2 & 1 & 3 \\
-\widetilde\xi_5 & 2 & 4 & 6 & 5 & 1 & 3 \\
-\widetilde\xi_6 & 2 & 4 & 6 & 5 & 4 & 3 \\ \hline
-\widetilde\xi'_1 & -4 & -5 & -6 & -4 & -2 & -3
\end{array}
\]
et $-\widetilde\xi'_i = -\widetilde\xi_i - 3r_0$, $3r_0 = (3, 6, 9, 6, 3,
6)$\footnote{Page 90, première version, les coefficients de $r_2$ sont repassés
(4, 4, 4, 4, 4, 9) ; la table de la page 92 les corrige, et on l'a vérifiée.
Pour les $\xi'_i$, la page écrit $-\widetilde\xi'_i = -\widetilde\xi_i - 3r_0
= -\widetilde\xi_i - \widetilde\eta$, avec $\widetilde\eta = 3(\eta - \lambda)
= -3r_0$ ; la seconde égalité a le mauvais signe, et la table des
$-\widetilde\xi'_i$ ($2, 7, 12, 8, 4, 9$, etc.) ajoute $3r_0$ au lieu de le
retrancher. On donne la première ligne corrigée ; les autres s'en déduisent.}.

La page 91, barrée mais lisible, retrouve les 27 comme les solutions de $x
\cdot \lambda = 1$, $x^2 = -1$ de la forme $a\eta + b\xi \pm \xi_i (\pm \xi_j)$
: (1) $x = a\eta + b\xi + \xi_i$ donne $\xi_i$ et $\xi'_i$ ; (3) $x = a\eta +
b\xi - \xi_i - \xi_j$ donne $\xi_{ij}$ ; (2) et (4) se ramènent aux précédents
par la symétrie $\delta \mapsto \tfrac23 \lambda - \delta$, qui est $-1$ sur
$E$ autour de $\lambda/3$ et préserve les deux équations. Il note encore une
application $r_i \mapsto -r_{6-i}$ ($i \leq 5$), $r_6 \mapsto r_{456}$ : on
vérifie que c'est un automorphisme du système de racines.

\pagerange{95}{98}
\section*{21. Les poids minuscules (pages 95 à 98)}

C'est la suite de la question de la page 57. Soit $R \subset E$ un système de
racines, $W$ son groupe de Weyl, $(\ ,\ )$ une forme invariante. On pose\footnote{Dans la définition de $\mathfrak{S}$ la page biffe le facteur
$2/(r,r)$, après avoir supposé $(r, r) = 2$ ; elle le réécrit dans la preuve
de $\mathfrak{S}_0 \subset \mathfrak{S}$, qui en a besoin. On le garde, ce qui
fait de $\mathfrak{S}$ l'ensemble des poids minuscules non nuls.}
\[
\begin{gathered}
\mathfrak{S} = \Bigl\lbrace x \in E \smallsetminus \lbrace 0\rbrace \Bigm|
\frac{2(x, r)}{(r, r)} \in \lbrace -1, 0, 1\rbrace \ \forall r \in R\Bigr\rbrace,
\\ \mathfrak{S}_0 = \lbrace x \neq 0 \mid wx - x \in R \text{ dès que } wx
\neq x\rbrace .
\end{gathered}
\]

\textbf{Faits généraux.} $\mathfrak{S}_0 \subset \mathfrak{S}$, puisque $s_r x
- x = -\tfrac{2(x, r)}{(r, r)}\, r$ doit être une racine, donc $\pm r$, dès
qu'il est non nul. $\operatorname{card} \mathfrak{S} \leq 3^{\operatorname{rg}}
- 1$ ; $\mathfrak{S}$ et $\mathfrak{S}_0$ sont stables par
$\operatorname{Aut}(E, R)$ et formés d'un nombre fini d'orbites. En langage
d'aujourd'hui, $\mathfrak{S}$ est la réunion des orbites des poids
fondamentaux minuscules\footnote{Le nom est de nous ; on ne sait pas s'il le
connaissait. Il ne figure pas sur ces pages.}.

\textbf{Cas $A_r$} (pages 95 et 96). Dans $E = \lbrace x \in \mathbb{R}^{r+1}
\mid \sum x_i = 0\rbrace$, un $x \in \mathfrak{S}$ prend deux valeurs
$-r_2/(r+1)$ ($r_1$ fois) et $r_1/(r+1)$ ($r_2$ fois), $r_1 + r_2 = r + 1$ :
ce sont les $r$ orbites des poids fondamentaux. Pour qu'il soit dans
$\mathfrak{S}_0$, il faut que $\sigma x - x$ soit une racine pour toute
permutation qui le bouge, c'est-à-dire $\operatorname{card}(I_1 \cap \sigma
I_2) = 1$ dès qu'il est non nul : d'où $r_1 = 1$ ou $r_2 = 1$. Pour $r \geq 2$,
$\mathfrak{S}_0$ a donc deux orbites, celles de la représentation standard et
de sa duale, échangées par $-1$.

\textbf{Cas $B_r$, $C_r$, $D_r$} (page 96). Pour $D_r$, $\mathfrak{S}$ a trois
orbites (les $\pm \varepsilon_i$ et les deux orbites de vecteurs $(\pm\tfrac12,
\ldots, \pm\tfrac12)$) et $\mathfrak{S}_0 = \emptyset$ pour $r \geq 4$, comme le
dit la page. Pour $B_r$ et $C_r$ la page conclut $\mathfrak{S} = \emptyset$
(pour $B_r$) et $\mathfrak{S}_0 = \emptyset$ (pour les deux) ; ce n'est vrai
que pour $\mathfrak{S}_0$ dans le cas $B_r$, $r \geq 3$. Pour $B_r$,
$\mathfrak{S}$ est l'orbite de $(\tfrac12, \ldots, \tfrac12)$ ; pour $C_r$,
$\mathfrak{S}$ est l'ensemble des $\pm\varepsilon_i$, et il est contenu dans
$\mathfrak{S}_0$, les différences $\pm\varepsilon_j - \varepsilon_i$ et
$-2\varepsilon_i$ étant des racines de $C_r$\footnote{La page, pour $B_r$,
écrit « $\xi_i$ égale 0 ou 2 » puis $\xi = (2, \ldots, 2)$, ce que la
transcription signale comme incohérent ; pour $C_r$ elle trouve $\lbrace \pm
\tfrac12\rbrace^r$, avec une normalisation différente. Les énoncés qu'on donne
sont les nôtres.}.

\textbf{Une variante} (pages 97 et 98). On affaiblit la condition :
\[
\mathfrak{S}'_0 = \lbrace x \mid w(\mathbb{R}x) \neq \mathbb{R}x \Rightarrow
wx - x \in R\rbrace \subset \mathfrak{S}' = \Bigl\lbrace x \Bigm| (x, r) \in
\tfrac{(r, r)}{2} \lbrace -1, 0, 1\rbrace \ \forall r \in R \smallsetminus
\mathbb{R}x\Bigr\rbrace ,
\]
et $\mathfrak{S}' = \mathfrak{S} \cup \lbrace \lambda s\rbrace$, où $s$ est une
racine et $\lambda > 0$ tel que $(\lambda s, r)$ soit dans $\tfrac{(r,r)}{2}
\lbrace -1, 0, 1\rbrace$ pour les racines $r \neq \pm s$. Un tel $\lambda$
existe si et seulement si les rapports $|(s, r)|/(r, r)$ sont égaux pour les
racines non orthogonales ni proportionnelles à $s$, c'est-à-dire si toutes les
racines ont même longueur ($A_r$, $D_r$, $E_6$, $E_7$, $E_8$), et alors
$\lambda = 1$ : dans le cas simplement lacé, $\mathfrak{S}' = \mathfrak{S}
\cup R$\footnote{La page 98 écrit « $\mathfrak{S}'_0 = \mathfrak{S}_0 \cup R$
si toutes les racines sont de même longueur ». Ce que la page 97 démontre est
l'énoncé sur $\mathfrak{S}'$, et c'est celui qu'on donne : une racine n'est
en général pas dans $\mathfrak{S}'_0$ (dans $A_2$, $s$ et une racine $s'$ avec
$(s, s') = -1$ donnent $s' - s$, qui n'est pas une racine).}. Pour $A_r$, enfin,
$\mathfrak{S}'_0 = \mathfrak{S}_0$ sauf pour $r = 3$, où s'ajoute l'orbite
$W \cdot (-\tfrac12, -\tfrac12, \tfrac12, \tfrac12)$ (le poids $\omega_2$,
dont l'opposé est dans l'orbite) ; on l'a vérifié.

\textbf{Pour $E_6$} --- que ces pages n'atteignent pas ---, $\mathfrak{S}$ a
exactement 54 éléments, les $\pm(\delta - \lambda/3)$, $\delta \in C$ : c'est
$\tfrac13 \widetilde\Delta$, la réunion des deux orbites de la page 57. La
question posée là a donc pour réponse : les 27 droites sont, à l'échelle $3$
et à translation par $\lambda$ près, l'une des deux orbites de poids minuscules
de $E_6$\footnote{Ce rapprochement est de nous, vérifié par machine ; les
pages ne le font pas. Noter que $\mathfrak{S}_0(E_6) = \emptyset$ : pour
deux droites qui se rencontrent, la différence n'est pas une racine. C'est
peut-être ce qui a conduit à la variante $\mathfrak{S}'_0$, sans plus de
succès.}.

\pagerange{100}{103}
\section*{22. Trois sous-groupes distingués, et un tripode (pages 100 à 103)}

\textbf{Page 100.} Soit $G$ un groupe et $N_1, N_2, N_3$ trois sous-groupes
distingués deux à deux d'intersection triviale et engendrant deux à deux $G$.
Alors $N_1 \times N_2 \to G$, $(x_1, x_2) \mapsto x_1 x_2$, est un isomorphisme,
$N_1 \simeq G/N_2$, $N_2 \simeq G/N_1$, et de même pour les autres
paires\footnote{La page écrit $G \simeq N_2 \rtimes N_1 \simeq N_3 \rtimes
N_1$, lecture incertaine du signe ; les produits sont directs, deux
sous-groupes distingués d'intersection triviale commutant entre eux.}. On peut
ajouter, ce que la page ne dit pas : $N_3$ est alors le graphe d'un
isomorphisme $N_1 \to N_2$, et sa normalité force $N_1$ à être commutatif ;
$G \simeq A \times A$ avec $A$ abélien.

\textbf{Page 101.} C'est ce que montre l'exemple $G \times G$ avec $N_2 = G
\times e$, $N_3 = e \times G$, $N_1 = \delta G$ la diagonale : $(g, h)(k, k)(g,
h)^{-1} = (gkg^{-1}, hkh^{-1})$, et la diagonale n'est distinguée que si $G$ est
commutatif. Suivent des essais avec un 1-cocycle $c : N_2 \to N_1$.

\textbf{Pages 102 et 103 : le tripode.} Trois ensembles $P, Q, R$ et une partie
$\Gamma \subset P \times Q \times R$ avec ses trois projections ; un point
$\gamma$ de $\Gamma$ définit des relations « $y \sim z$ le long de $p(\gamma)$ »,
et la page demande si deux telles relations en entraînent une troisième, puis
fait le calcul pour $\Gamma = \lbrace abc = 1\rbrace$. Un diagramme reprend
$\widetilde\Gamma/\sigma \to \widetilde P_u/\sigma$, \ldots : c'est la
situation du lemme 7. La page 103 écrit une configuration de huit triplets
dans $P \times Q \times R$ et la question : « pour tous $a_0, \bar r, \bar r'
\in P$, $b_0, \tilde r'' \in Q$, existe-t-il $\bar\pi \in P$, $\tilde r',
\tilde\sigma \in Q$, $r', \pi, \sigma, \rho \in R$ tels que l'on ait » ces
triplets. C'est une condition de fermeture de type Reidemeister, qui
caractérise les carrés latins provenant d'un groupe\footnote{Le nom et
l'identification sont de nous, et la configuration de la page, surchargée et
en partie hachurée, n'a pas été comparée terme à terme avec celle de
Reidemeister (1929) ; la condition de Thomsen caractérise de même les groupes
commutatifs. Références citées de mémoire, ainsi que le livre de Blaschke et
Bol sur les tissus (1938).}.

\pagerange{105}{105}
\section*{23. Un aparté : le $56$ de $E_7$ (page 105)}

$\dim E_7 = 133$, avec 126 racines. 56 indéterminées $x_i, y_i$ ($1 \leq i
\leq 7$), $x_{ih} = -x_{hi}$, $y_{ih} = -y_{hi}$ ($1 \leq i < h \leq 7$) :
$7 + 7 + 21 + 21 = 56$, la représentation de dimension 56 de $E_7$ décomposée
sous un sous-groupe de type $A_6$. Il écrit un invariant quartique,
\[
\sum x_i y_j x_{ih} y_{jk} + \tfrac14 \sum x_{\lambda\mu} x_{\nu\rho}
y_{\mu\nu} y_{\lambda\rho} + \sum \bigl(x_i y_{\lambda\mu} y_{\nu\rho}
y_{\sigma\tau} + y_i x_{\lambda\mu} x_{\nu\rho} x_{\sigma\tau}\bigr),
\]
la dernière somme portant sur les permutations paires $(i, \lambda, \mu, \nu,
\rho, \sigma, \tau)$ de $[1, 7]$, et la forme symplectique $\sum (x_i\, dy_i -
y_i\, dx_i) + \sum (x_{ih}\, dy_{ih} - y_{ih}\, dx_{ih})$\footnote{Les indices
du premier terme sont incertains ($y_{jk}$ ou $y_{jh}$). Ce sont l'invariant
quartique de Cartan--Freudenthal et la forme symplectique invariante sur la
représentation de dimension 56 ; on n'a vérifié ni les coefficients ni
l'invariance. Le rapport avec les 27 droites est celui de $E_6 \subset E_7$,
la page ne le dit pas.}.

\pagerange{107}{118}
\section*{24. Trois torseurs et une correspondance (pages 107 à 118)}

\subsection*{Le cas commutatif (pages 107 à 109)}

Soit $G$ un groupe commutatif, $P_1, P_2, P_3$ trois $G$-torseurs, et $\Gamma
\subset P_1 \times P_2 \times P_3$ un sous-torseur sous $\operatorname{Ker}(G^3
\xrightarrow{+} G)$ --- ce qui revient à une trivialisation $t$ du produit de
Baer $P_1 \wedge P_2 \wedge P_3$. Il annonce que le foncteur
\[
(G, P_1, P_2, P_3, t) \longmapsto (P_1, P_2, P_3, \Gamma)
\]
vers les quadruplets de trois ensembles et d'une partie de leur produit est
pleinement fidèle. L'image vérifie la condition (1) : les projections $p_{ij} :
\Gamma \to P_i \times P_j$ sont bijectives ; se donner $\Gamma$ vérifiant (1)
revient à se donner $\pi_3 : P_1 \times P_2 \to P_3$ dont toutes les
applications partielles sont bijectives, d'où $h_{23} : P_1 \to
\operatorname{Isom}(P_2, P_3)$, etc. Autrement dit, $\Gamma$ est un \emph{carré
latin} (la table d'un quasigroupe), et la question est de reconnaître ceux qui
viennent d'un groupe\footnote{Noms de nous. Le cas du lemme 7 est celui de $G
= V = \mathbb{F}_2^2$ : $\Gamma \subset P_u \times P_v \times P_w$ y est la
table du groupe de Klein.}.

Pour un groupe $G$ quelconque, on se ramène (page 109) aux trois torseurs
triviaux et à
\[
\Gamma_0 = \lbrace (x, y, z) \in G^3 \mid xyz = 1\rbrace = \lbrace (g_2
g_3^{-1}, g_3 g_1^{-1}, g_1 g_2^{-1}) \mid g_i \in G\rbrace ,
\]
l'orbite de $(1, 1, 1)$ sous l'action $\sigma(g_1, g_2, g_3)(x_1, x_2, x_3) =
(g_2 x_1 g_3^{-1}, g_3 x_2 g_1^{-1}, g_1 x_3 g_2^{-1})$\footnote{La page
appelle « condition (A) » le fait que $\Gamma$ soit ainsi « linéaire » ; la
condition (A) n'est énoncée nulle part sur ces feuillets.}. Un point $\gamma =
(a_1, a_2, a_3)$ de $\Gamma$ définit des isomorphismes $\tilde a_1 : G_2 \to
G_3$, \ldots, par $g_2 a_1 = a_1 g_3$ ; si on remplace $\gamma$ par
$\sigma(\alpha)(\gamma)$, ils deviennent
\[
\tilde a'_1 = \operatorname{int}(\alpha_3)\, \tilde a_1\,
\operatorname{int}(\alpha_2)^{-1}, \qquad \text{et circulairement},
\]
calcul refait page 116.

\subsection*{Les autotopies (pages 108, 116 et 118)}

Le groupe $\operatorname{Aut}(G) \ltimes G^3$, de loi
\[
(u, \alpha, \beta, \gamma)(u', \alpha', \beta', \gamma') = (uu', \alpha
u(\alpha'), \beta u(\beta'), \gamma u(\gamma')),
\]
d'inverse $(u^{-1}, u^{-1}(\alpha^{-1}), u^{-1}(\beta^{-1}), u^{-1}(\gamma^{-1}))$
(page 118, qui vérifie l'associativité), opère sur $G$ en trois exemplaires
par
\[
(u, \alpha, \beta, \gamma) \longmapsto \bigl(x \mapsto \beta u(x)
\gamma^{-1},\ x \mapsto \gamma u(x) \alpha^{-1},\ x \mapsto \alpha u(x)
\beta^{-1}\bigr) ,
\]
en préservant $\Gamma_0$, et le noyau est $N = \lbrace (\operatorname{int}(g),
g^{-1}, g^{-1}, g^{-1}) \mid g \in G\rbrace$ (page 108)\footnote{Les indices du
premier composant sont repassés ; on les donne selon la permutation circulaire
des deux autres, comme la transcription. Les quotients par $N$ sont notés
$\Gamma(G)$ (avec $\operatorname{Int} G$) et $\widehat\Gamma(G)$ (avec
$\operatorname{Aut} G$) ; le second dénominateur est laissé vide.}. La page
116 montre que ce sont toutes les \emph{autotopies} de $\Gamma_0$ : si $(\varphi,
\psi, \chi)$ sont trois bijections de $G$ telles que $xyz = 1 \Rightarrow
\varphi(x) \psi(y) \chi(z) = 1$, on les normalise par des translations
pour que $\varphi'(1) = \psi'(1) = \chi'(1) = 1$ ; alors $\psi'(y) =
\chi'(y^{-1})^{-1}$, de même pour les autres, d'où $\varphi' = \psi' = \chi' =
u$, et $u(x) u(y) = u(xy)$. Le groupe des autotopies de la table de $G$ est
donc $(\operatorname{Aut} G \ltimes G^3)/N$, d'ordre $|\operatorname{Aut}
G| \cdot |G|^2$\footnote{C'est un théorème classique sur les quasigroupes
(on le rattache en général à Albert, 1943, cité de mémoire) ; le nom
d'« autotopie » est de nous.}.

\subsection*{Brouillons (pages 110, 111, 114, 115)}

Page 110 (tête-bêche) : l'opération $x \circ y = -x - y$ sur un groupe
abélien, commutative et non associative ($(x \circ y) \circ z = x + y - z$),
et une opération $x \ast y = \overline{x \circ y}$ qui rend la loi de groupe ;
les tables de $\mathbb{Z}/4$ et (page 111) de $\mathbb{Z}/5$. Page 111 : le
triangle de sommets $P_1, P_2, P_3$ et de côtés $G_1, G_2, G_3$, et
$P_1 \wedge_{G_3} P_2 \wedge_{G_1} P_3 \simeq \mathbf{1}_{G_2}$. Pages 114 et
115 : pour un carré latin d'ordre $n$ et un point base $s = (a_1, a_2, a_3)$,
$\Gamma = \lbrace s\rbrace \cup \Gamma_1 \cup \Gamma_2 \cup \Gamma_3 \cup
\Gamma^*$, avec $\Gamma_i = p_i^{-1}(a_i) \smallsetminus \lbrace s\rbrace$
($n - 1$ éléments chacun) et $\Gamma^*$ le reste ($n^2 - 3n + 2$) ; il attache
à chaque point un objet $E_x$ avec des isomorphismes le long de $\Gamma$,
normalisés en $s$, et compte $3(n - 1) + 2(n^2 - 3n + 2) = 2n^2 - 3n + 1$
isomorphismes libres\footnote{L'interprétation de ce compte est de nous ;
pour $n = 4$, le cas des $P_u$, il vaut 21, nombre écrit à côté (« 21/42 »).}.

\subsection*{Bitorseurs (page 112)}

L'énoncé propre, pour des groupes quelconques : $P_1$ est un bitorseur sous
$(G_2, G_3)$, $P_2$ sous $(G_3, G_1)$, $P_3$ sous $(G_1, G_2)$ ; les produits
contractés circulaires $P_1 \wedge P_2 \wedge P_3$ (bitorseur sous $(G_2,
G_2)$), $P_2 \wedge P_3 \wedge P_1$, $P_3 \wedge P_1 \wedge P_2$ sont munis de
trivialisations $\varphi_2, \varphi_3, \varphi_1$, données équivalentes ; et,
en termes ensemblistes, $\Gamma \subset P_1 \times P_2 \times P_3$ est une
orbite non vide de $G_1 \times G_2 \times G_3$ opérant par $\sigma$ comme
ci-dessus. Un point $\gamma$ de $\Gamma$ fournit les isomorphismes $\tilde
a_1 : G_2 \to G_3$, $\tilde a_2 : G_3 \to G_1$, $\tilde a_3 : G_1 \to G_2$, et
leur changement quand $\gamma$ varie est donné par la formule précédente
(« opér.\ transitive d'iso. »)\footnote{Les bitorseurs sont ceux de la
cohomologie non abélienne de Giraud (1971), que la page n'invoque pas ;
l'opposé $\mathring{P}_1$ d'un bitorseur et les produits dans l'ordre inverse
sont écrits à droite de la page, dans un ordre des facteurs qui ne suit pas
une règle unique.}.

\subsection*{Retour au graphe (pages 113 et 117)}

La page 113 referme la boucle. Soit $V$ de dimension 2 sur $\mathbb{F}_2$, $I$
un ensemble à trois éléments, $(P_i)_{i \in I}$ des $V$-torseurs avec une
trivialisation de $\bigwedge P_i$, et $\Gamma = \lbrace (x_i) \mid \sum x_i =
0\rbrace$ (16 éléments) ; des revêtements doubles $\widetilde P_i \to P_i$,
$\widetilde\Gamma \to \Gamma$, avec $\widetilde\Gamma \simeq \Gamma \times_{P_i}
\widetilde P_i$ pour chaque $i$ ; sur $\widetilde P = \coprod \widetilde P_i$ (24
points), $x \neq y$ sont liés s'ils sont conjugués dans un même $\widetilde
P_i$, ou images d'un même point de $\widetilde\Gamma$ dans deux $\widetilde
P_i$ différents. \textbf{Axiome} : pour $i \neq j$, l'application $\alpha_{ji}$
qui associe à $\tilde x \in \widetilde P_i$ la section de $\widetilde P_j \to
P_j$ formée de ses voisins est injective, son image est un torseur sous le
sous-groupe $\operatorname{Ker}(\mathbb{F}_2^{P_j} \xrightarrow{\Sigma}
\mathbb{F}_2)$ des fonctions paires, et les images de $\widetilde P_i$ et de
$\widetilde P_k$ sont disjointes. C'est exactement ce que démontraient les
lemmes 3 et 4 des pages 2 à 4 : $\alpha_{ji}(x) = \omega_j(x)$, deux sections
de la même image diffèrent sur 2 ou 4 paires, et les 16 sections se
partagent en deux classes de parité de 8, l'une pour $\widetilde P_i$, l'autre
pour $\widetilde P_k$\footnote{Ce rapprochement est de nous ; on l'a vérifié
sur le modèle.}. La page 117 redessine le triangle $u, v, w$ avec ses quatre
paires d'arêtes $(\alpha_i, \alpha'_i)$, $(\beta_i, \beta'_i)$, $(\gamma_i,
\gamma'_i)$ à chaque sommet, recopie les tables (I), (I$'$) et la ligne
« $\alpha_1$ lié à $\beta'_1, \beta_2, \beta_3, \beta_4$ » des pages 9 et 10,
et reprend le calcul de $\tilde a'_1$.

\pagerange{119}{119}
\section*{25. Un autre sujet (page 119)}

Tête-bêche, au crayon et à l'encre, des croquis et des formules sur un autre
sujet : $S^2 \times S^2$ privé de la diagonale et de l'antidiagonale, fibré
au-dessus de $S^2$ de fibre $S^2$ moins deux points, $\simeq \mathbb{C}^*$, et
le début de la suite exacte des groupes fondamentaux ; les idempotents $p$ de
rang 1 ($p^2 = p$, $\operatorname{Tr} p = 1$), qui correspondent aux
décompositions $E \simeq L \oplus L'$ d'un plan en deux droites ; une suite
$0 \to F \to E_p \to \mathcal{O}(1) \to 0$ avec $F \simeq \mathcal{O}(-1)$. La
page ne conclut rien\footnote{L'espace des idempotents de rang 1 de
$M_2(\mathbb{C})$ est $\mathbb{P}^1 \times \mathbb{P}^1$ privé de la diagonale
(image et noyau) ; le lien avec le $S^2 \times S^2$ de la page est notre
lecture.}.

\end{document}
